%

\chapter{Konstrukcje klasyczne} % lang-pl
\chapter{Costruzioni classiche} % lang-it

% TODO: https://en.wikipedia.org/wiki/Straightedge_and_compass_construction
% https://en.wikipedia.org/wiki/Compass_equivalence_theorem
% https://en.wikipedia.org/wiki/4,294,967,295
Konstrukcje klasyczne zwane też platońskimi to wspólna nazwa problemów, których celem jest wyznaczenie odcinków, kątów czy też okręgów spełniających dane warunki, korzystając przy tym jedynie z linijki (bez podziałki!) oraz cyrkla. % lang-pl
Le costruzioni classiche, dette anche platoniche, sono il nome collettivo di problemi il cui scopo è determinare segmenti, angoli oppure cerchi che soddisfano determinate condizioni utilizzando esclusivamente una riga (non graduata!) e un compasso. % lang-it

Współczesny cyrkiel jest sztywny (po angielsku: \emph{rigid compass, modern compass}) i potrafi przenosić odcinki, bo zachowuje swoje rozwarcie po podniesieniu z płaszczyzny rysunku. % lang-pl
Starożytni skorzystają z cyrkla zapadającego się (po angielsku: \emph{traditional collapsing compass, divider}), ale nie ograniczy to ich możliwości, ponieważ Euklides udowodni bardzo szybko twierdzenie o równoważności obydwu cyrkli (I.2). % lang-pl
Dowód przedstawia też Eves \cite[s. 155]{eves1_1972}. % lang-pl
% rigid compass = modern compass
% traditional collapsing compass = divider
Il compasso moderno è rigido (\emph{rigid compass, modern compass}) e permette di trasportare segmenti, poiché mantiene la propria apertura anche dopo essere stato sollevato dal piano del disegno. % lang-it
Gli antichi utilizzavano invece il compasso collassabile (\emph{traditional collapsing compass, divider}), ma ciò non limiterà le loro possibilità, poiché Euclide dimostrerà molto presto il teorema dell'equivalenza dei due tipi di compasso (I.2). % lang-it
Una dimostrazione è presentata anche da Eves \cite[p. 155]{eves1_1972}. % lang-it

Obydwa przyrządy są wyidealizowane: linijka jest nieskończenie długa, ma tylko jeden brzeg i~jak już napisaliśmy, nie ma podziałki. % lang-pl
Cyrkiel natomiast może kreślić okręgi o dowolnie małych lub dużych promieniach. % lang-pl
Zabronione jest korzystanie jednocześnie z cyrkla i linijki, ponieważ razem są w stanie imitować linijkę z podziałką. % lang-pl
Ślad, jaki zostawiają na kartce, jest nieskończenie cienki. % lang-pl
Entrambi gli strumenti sono idealizzati: la riga è infinitamente lunga, possiede un solo bordo e, come già detto, non è graduata. % lang-it
Il compasso, invece, può tracciare cerchi di raggio arbitrariamente piccolo o grande. % lang-it
Non è consentito utilizzare contemporaneamente riga e compasso, poiché insieme sono in grado di imitare una riga graduata. % lang-it
La traccia che lasciano sul foglio è infinitamente sottile. % lang-it

% eksperymentalne tłumaczenie na włoski

%

\section{Łatwe} % lang-pl
\section{Problemi facili} % lang-it


Hartshorne \cite[s. 122]{hartshorne2000} napisze, że Kartezjusz około 1637 pokazał, że: % lang-pl
Hartshorne \cite[p. 122]{hartshorne2000} scrive che Cartesio mostrò intorno al 1637 che: % lang-it

\begin{proposition}
    Niech $P_1 = (x_1, y_1)$, $P_2 = (x_2, y_2)$, \ldots, $P_n = (x_n, y_n)$ będą punktami na rzeczywistej płaszczyźnie kartezjańskiej, na której dane są również punkty $(0, 0)$ oraz $(1, 0)$. % lang-pl
    
    Wtedy następujące warunki są równoważne: % lang-pl
    Siano $P_1 = (x_1, y_1)$, $P_2 = (x_2, y_2)$, \ldots, $P_n = (x_n, y_n)$ punti del piano cartesiano reale, sul quale sono inoltre assegnati i punti $(0,0)$ e $(1,0)$. % lang-it
    
    Allora le seguenti condizioni sono equivalenti: % lang-it
    \begin{itemize}
        \item punkt $P = (x, y)$ można wykreślić przy pomocy cyrkla i linijki; % lang-pl
        \item liczby $x, y$ można otrzymać z $a_1, a_2, \ldots, a_n, b_1, b_2, \ldots, b_n$ przy pomocy dodawania, odejmowania, mnożenia, dzielenia lub rozwiązując skończoną liczbę liniowych lub kwadratowych równań, w których wyciągamy pierwiastki kwadratowe jedynie z dodatnich liczb. % lang-pl
        \item il punto $P=(x,y)$ può essere costruito con riga e compasso; % lang-it
        \item i numeri $x,y$ possono essere ottenuti da $a_1, a_2, \ldots, a_n, b_1, b_2, \ldots, b_n$ mediante addizione, sottrazione, moltiplicazione, divisione oppure risolvendo un numero finito di equazioni lineari o quadratiche, nelle quali si estraggono radici quadrate soltanto di numeri positivi. % lang-it
    \end{itemize}
\end{proposition}

% TODO: Hartshorne s. 103
\begin{geoconstruction}
    Dany jest kąt. % lang-pl
    
    Znaleźć jego dwusieczną. % lang-pl
    È dato un angolo. % lang-it
    
    Costruirne la bisettrice. % lang-it
\index{dwusieczna}%
\end{geoconstruction}

Hartshorne \cite[s. 23]{hartshorne2000}. % lang-pl

Hartshorne \cite[p. 23]{hartshorne2000}. % lang-it
% TODO: Zetel s. 21 to samo, ale wierzchołek jest niedostępny

\begin{problem}
    Środek odcinka. % lang-pl
    
    Il punto medio di un segmento. % lang-it
\end{problem}

Hartshorne \cite[s. 23]{hartshorne2000}. % lang-pl

Hartshorne \cite[p. 23]{hartshorne2000}. % lang-it

\begin{problem}
    Prosta prostopadła do prostej, przechodząca przez punkt. % lang-pl
    
    La retta perpendicolare a una retta assegnata passante per un punto. % lang-it
    \index{prosta!prostopadła}
\end{problem}

Hartshorne \cite[s. 23, 24]{hartshorne2000}. % lang-pl

Hartshorne \cite[p. 23, 24]{hartshorne2000}. % lang-it

\begin{problem}
    Prosta równoległa do prostej, przechodząca przez punkt. % lang-pl
    
    La retta parallela a una retta assegnata passante per un punto. % lang-it
    \index{prosta!równoległa}
\end{problem}

Hartshorne \cite[s. 24]{hartshorne2000}. % lang-pl
Hartshorne \cite[p. 24]{hartshorne2000}. % lang-it
\begin{problem} % lang-it
    Il centro di un cerchio assegnato. % lang-it
\end{problem} % lang-it


\begin{problem} % lang-pl
    Środek danego okręgu. % lang-pl
\end{problem} % lang-pl
Hartshorne \cite[p. 24]{hartshorne2000}. % lang-it

Hartshorne \cite[s. 24]{hartshorne2000}. % lang-pl
Porównaj z zadaniem Napoleona \ref{napoleon_problem}. % lang-pl
Confrontare con il problema di Napoleone \ref{napoleon_problem}. % lang-it
\index{zadanie!Napoleona}

\begin{problem}
    Styczna do okręgu przez punkt poza nim. % lang-pl
    
    La tangente a un cerchio passante per un punto esterno. % lang-it
    \index{styczna}
\end{problem}

Hartshorne \cite[s. 24]{hartshorne2000}. % lang-pl
Hartshorne \cite[p. 24]{hartshorne2000}. % lang-it
\begin{problem} % lang-it
    Sono dati due punti $A$, $B$ e una retta $k$. % lang-it
    Trovare un cerchio tangente alla retta e passante per entrambi i punti. % lang-it
\end{problem} % lang-it


Triangolo inscritto: 13 passi, circoscritto: 7 passi. % lang-it
\begin{problem}
    Dane są dwa punkty $A$, $B$ i prosta $k$. % lang-pl
    
    Znaleźć okrąg styczny do prostej, który przechodzi przez obydwa punkty. % lang-pl
    È dato un triangolo con un lato distinto. % lang-it
    
    Inscrivere in esso un quadrato in modo che uno dei suoi lati giaccia sul lato distinto del triangolo. % lang-it
\end{problem}

Trójkąt wpisany 13, opisany 7. % lang-pl
Eves \cite[p. 162]{eves1_1972} % lang-it
\begin{problem} % lang-it
    Sono dati due triangoli, uno contenuto nell'altro. % lang-it
    Inscrivere nel triangolo esterno un terzo triangolo circoscritto a quello interno. % lang-it
\end{problem} % lang-it


Eves \cite[p. 167]{eves1_1972} % lang-it
\begin{problem}
    Dany jest trójkąt z wyróżnionym bokiem. % lang-pl
    
    Wpisać w niego kwadrat tak, by jeden z jego czterech boków leżał na wyróżnionym boku trójkąta. % lang-pl
    Sono dati una retta, un segmento e un punto esterno alla retta. % lang-it
    
    Costruire un triangolo isoscele con vertice nel punto assegnato e base, situata sulla retta, uguale al segmento dato. % lang-it
\end{problem}

Eves \cite[s. 162]{eves1_1972} % lang-pl
Hartshorne \cite[p. 25]{hartshorne2000}. % lang-it
\begin{problem} % lang-it
    Sono dati una retta passante per $B$ e un punto $A$ esterno ad essa. % lang-it
    Trovare un cerchio passante per entrambi i punti e tangente alla retta. % lang-it
\end{problem} % lang-it


Hartshorne \cite[p. 25]{hartshorne2000}. % lang-it
\begin{problem}
    Dane są dwa trójkąty, jeden zawarty w drugim. % lang-pl
    
    Wpisać w zewnętrzny trójkąt trzeci trójkąt tak, by był opisany na wewnętrznym trójkącie. % lang-pl
    Tre cerchi che si intersecano a coppie ortogonalmente. % lang-it
\end{problem}

Eves \cite[s. 167]{eves1_1972} % lang-pl

Problema di Hartshorne \cite[p. 25]{hartshorne2000}; una soluzione di Łukasz Rajkowski in $\Delta_{24}^{12}$ utilizza un fascio di cerchi. % lang-it

\begin{problem}
    Dana jest prosta, odcinek i punkt poza prostą. % lang-pl
    
    Wykreślić trójkąt równoramienny, z wierzchołkiem w punkie i z podstawą równą odcinkowi na prostej. % lang-pl
\end{problem} % lang-pl
Hartshorne \cite[s. 25]{hartshorne2000}. % lang-pl
\begin{problem} % lang-pl
    Dana jest prosta przez punkt $B$, punkt $A$ poza prostą. % lang-pl
    Znaleźć okrąg, który przechodzi przez obydwa punkty, styczny do prostej. % lang-pl
\end{problem} % lang-pl
Hartshorne \cite[s. 25]{hartshorne2000}. % lang-pl
\begin{problem} % lang-pl
    Trzy okręgi przecinające się parami pod kątem prostym. % lang-pl
\end{problem} % lang-pl
Problem Hartshorne'a \cite[s. 25]{hartshorne2000}, rozwiązanie Łukasza Rajkowskiego $\Delta_{24}^{12}$ z wykorzystaniem pęku okręgów. % lang-pl
\begin{problem} % lang-pl
    Podzielić odcinek na trzy równe części. % lang-pl
    Dividere un segmento in tre parti uguali. % lang-it
    \index{trysekcja!odcinka}
\end{problem}

Hartshorne \cite[s. 25]{hartshorne2000}. % lang-pl

Hartshorne \cite[p. 25]{hartshorne2000}. % lang-it

% TODO: Neugebauer s. 67
\begin{problem}
    Okrąg styczny do prostej, przechodzący przez dwa punkty poza nią. % lang-pl
    
    Un cerchio tangente a una retta e passante per due punti esterni ad essa. % lang-it
\end{problem}

Hartshorne \cite[s. 43]{hartshorne2000}. % lang-pl
Hartshorne \cite[p. 43]{hartshorne2000}. % lang-it
\begin{problem} % lang-it
    Un cerchio tangente a due rette e passante per un punto esterno ad esse. % lang-it
\end{problem} % lang-it


Hartshorne \cite[p. 44]{hartshorne2000}. % lang-it
\begin{problem}
    Okrąg styczny do dwóch prostych, przechodzący przez punkt poza nimi. % lang-pl
    
    Un rettangolo avente la stessa area di un triangolo assegnato. % lang-it
    È assegnato uno dei lati del rettangolo. % lang-it
\end{problem}

Hartshorne \cite[s. 44]{hartshorne2000}. % lang-pl
Hartshorne \cite[p. 43]{hartshorne2000}. % lang-it
\begin{problem} % lang-it
    Un quadrato avente la stessa area di un rettangolo assegnato. % lang-it
\end{problem} % lang-it


Hartshorne \cite[p. 43]{hartshorne2000}. % lang-it
\begin{problem}
    Prostokąt o polu takim, jak dany trójkąt. % lang-pl
    
    Dany jest jeden bok prostokąta. % lang-pl
    Dividere un triangolo mediante una retta passante per un punto assegnato in due parti di uguale area. % lang-it
\end{problem}

Hartshorne \cite[s. 43]{hartshorne2000}. % lang-pl

Hartshorne \cite[p. 44]{hartshorne2000}. % lang-it

\begin{problem}
    Kwadrat o polu takim, jak dany prostokąt. % lang-pl
    
\end{problem} % lang-pl
Hartshorne \cite[s. 43]{hartshorne2000}. % lang-pl
\begin{problem} % lang-pl
    Podzielić trójkąt prostą przez dany punkt na dwie części o równych polach. % lang-pl
    Sono dati un segmento $AB$ e un punto $P$ all'interno di un cerchio. % lang-it
    
\end{problem} % lang-pl
Hartshorne \cite[s. 44]{hartshorne2000}. % lang-pl
\begin{problem} % lang-pl
    Dany jest odcinek $AB$ oraz punkt $P$ wewnątrz okręgu. % lang-pl
    Skonstruować cięciwę tego okręgu, która przechodzi przez punkt $P$ o długości takiej samej, jak odcinek $AB$. % lang-pl
    Costruire una corda del cerchio passante per $P$ e avente la stessa lunghezza del segmento $AB$. % lang-it
    \index{cięciwa}
\end{problem}

\begin{problem}
    Dany jest odcinek $AB$, inny odcinek o długości $d$ oraz kąt $\alpha$. % lang-pl
    
    Skonstruować trójkąt $ABC$ tak, by kąt przy wierzchołku $C$ miał miarę $\alpha$, zaś suma długości ramion tego kąta była równa $d$. % lang-pl
    Sono dati un segmento $AB$, un altro segmento di lunghezza $d$ e un angolo $\alpha$. % lang-it
    
    Costruire un triangolo $ABC$ tale che l'angolo in $C$ abbia ampiezza $\alpha$ e che la somma delle lunghezze dei lati che formano tale angolo sia uguale a $d$. % lang-it
\end{problem}

\begin{problem}
    Dane są dwa okręgi takie, że żaden nie jest zawarty w drugim. % lang-pl
    
    Skonstruować styczną do obydwu okręgów. % lang-pl
    Sono dati due cerchi, nessuno dei quali è contenuto nell'altro. % lang-it
    
    Costruire una tangente comune ai due cerchi. % lang-it
    \index{dwustyczna}
\end{problem}

\begin{problem}
    Dany jest okrąg $\Gamma$ oraz jego środek $O$. % lang-pl
    
    Skonstruować trzy przystające okręgi, które są styczne do pozostałych dwóch oraz do $\Gamma$. % lang-pl
    Sono dati un cerchio $\Gamma$ e il suo centro $O$. % lang-it
    
    \hfill \emph{(13 kroków)}. % lang-pl
    Costruire tre cerchi congruenti tangenti tra loro a coppie e tangenti a $\Gamma$. % lang-it
    \hfill \emph{(13 passi)}. % lang-it
\end{problem}

\begin{problem}
    Dany jest okrąg $\Gamma$ oraz dwa punkty $A$ i $B$. % lang-pl
    
    Skonstrować punkt $C$ na okręgu $\Gamma$ tak, by odcinek łączący punkty przecięcia prostych $CA$, $CB$ z okręgiem $\Gamma$ był równoległy do odcinka $AB$. % lang-pl
    Sono dati un cerchio $\Gamma$ e due punti $A$ e $B$. % lang-it
    
    Costruire un punto $C$ sul cerchio $\Gamma$ tale che il segmento che unisce i secondi punti d'intersezione delle rette $CA$ e $CB$ con $\Gamma$ sia parallelo ad $AB$. % lang-it
\end{problem}

\begin{problem}
    Skonstruować trzy parami styczne okręgi, każdy o innym promieniu, których środki nie są współliniowe. % lang-pl
    
    \hfill \emph{(7 kroków)}. % lang-pl
    Costruire tre cerchi mutuamente tangenti, tutti di raggio diverso, con centri non allineati. % lang-it
    \hfill \emph{(7 passi)}. % lang-it
\end{problem}

Skonstruować oś potęgową. % lang-pl

Costruire l'asse radicale. % lang-it
\index{oś!potęgowa}
Audin \cite[s. 107]{audin_2003} podaje ten fakt w formie ćwiczenia. % lang-pl

Audin \cite[p. 107]{audin_2003} presenta questo fatto come esercizio. % lang-it

\begin{problem}
\label{konstrukcja_kwadratu_wpisanego}%
    Dany jest trójkąt $ABC$ z kątami ostrymi przy $B$ i $C$. % lang-pl
    
    Skonstruować kwadrat, którego jeden bok leży na odcinku $BC$, zaś pozostałe dwa wierzchołki leżą na odcinkach $AC$ i $AB$. % lang-pl
    È dato un triangolo $ABC$ con angoli acuti in $B$ e $C$. % lang-it
    
    Costruire un quadrato avente un lato sul segmento $BC$ e gli altri due vertici sui segmenti $AC$ e $AB$. % lang-it
\end{problem}

(Po tej konstrukcji warto przeczytać o trójkącie Calabiego z faktu \ref{trojkat_calabiego} raz jeszcze). % lang-pl

(Dopo questa costruzione vale la pena rileggere il fatto \ref{trojkat_calabiego} sul triangolo di Calabi). % lang-it

%
% eksperymentalne tłumaczenie na włoski

%

\section{Trudniejsze} % lang-pl
\section{Problemi più difficili} % lang-it
% eksperymentalne tłumaczenie na włoski

%

\subsection{Okręgi Apoloniusza i Soddy'ego} % lang-pl

\subsection{Cerchi di Apollonio e di Soddy} % lang-it

% https://en.wikipedia.org/wiki/Problem_of_Apollonius
\begin{problem}
    \label{problem_apoloniusza}%
    Na płaszczyźnie dane są trzy okręgi. % lang-pl
    Skonstruować czwarty, który jest styczny do każdego z~pozostałych. % lang-pl
    
    \index{okrąg!Apoloniusza}% % lang-pl
    Nel piano sono dati tre cerchi. % lang-it
    Costruire un quarto che sia tangente a ciascuno degli altri. % lang-it
    \index{cerchio!di Apollonio}% % lang-it
\end{problem}

Dyskusję na temat problemu znajdziemy u Guzickiego \cite[s. 444-461]{guzicki_2021}. % lang-pl
Eves \cite[s. 164]{eves1_1972} rozważa szczególny przypadek trzech współpękowych okręgów. % lang-pl
W ogólności okręgi mogą mieć zerowy promień (i skurczyć się do punktu) lub nieskończony promień (i wydłużyć się do prostej). % lang-pl
Zależnie od wyboru kształtów, mamy do rozwiązania dziesięć różnych zadań: PPP, PPL, PLL, LLL, PPC, PLC, LLC, PCC, LCC, CCC (klasyczne zadanie Apoloniusza). % lang-pl
Tutaj P oznacza punkt, L prostą, zaś C okrąg, od pierwszych liter stosownych słów po angielsku.% % lang-pl
Una discussione del problema si trova in Guzicki \cite[s. 444-461]{guzicki_2021}. % lang-it
Eves \cite[s. 164]{eves1_1972} considera il caso particolare di tre cerchi coassiali. % lang-it
In generale i cerchi possono avere raggio nullo (riducendosi a un punto) oppure raggio infinito (diventando una retta). % lang-it
A seconda della scelta delle figure, si ottengono dieci problemi diversi: PPP, PPL, PLL, LLL, PPC, PLC, LLC, PCC, LCC, CCC (il problema classico di Apollonio). % lang-it
Qui P indica un punto, L una retta e C un cerchio, dalle iniziali dei termini inglesi corrispondenti.% % lang-it

Z zadaniem wiąże się dużo późniejszy wynik. % lang-pl
W dwóch listach z 1643 roku do czeskiej księżniczki Elżbiety z Palatynatu Kartezjusz przedstawi uproszczoną wersję, w której trzy okręgi są parami styczne i znajdzie (bez dowodu) relację wiążącą promienie okręgów.% % lang-pl
\index[persons]{van Pallandt, Elżbieta (z Palatynatu)}% % lang-pl
Al problema è legato un risultato molto più tardo. % lang-it
In due lettere del 1643 alla principessa Elisabetta di Boemia, Cartesio presentò una versione semplificata in cui tre cerchi sono a due a due tangenti e trovò (senza dimostrazione) una relazione tra i loro raggi.% % lang-it
\index[persons]{di Boemia, Elisabetta (Principessa Palatina)}% % lang-it

\begin{theorem}[Kartezjusza] % lang-pl
Dane są trzy okręgi, parami styczne o promieniach $r_1, r_2, r_3$ oraz czwarty styczny do każdego z nich\footnote{istnieją dwa takie, jeden z~nich zazwyczaj styczny wewnętrznie, a w szczególnym przypadku może okazać się prostą.}. % lang-pl
Definiujemy krzywiznę ze znakiem jako $\pm 1 / r$, gdzie znak $-$ odpowiada okręgowi stycznemu zewnętrznie. % lang-pl
Wtedy% % lang-pl
\begin{theorem}[di Cartesio] % lang-it
Sono dati tre cerchi, a due a due tangenti, con raggi $r_1, r_2, r_3$ e un quarto tangente a ciascuno di essi\footnote{ne esistono due, uno dei quali è di solito tangente internamente e, in un caso particolare, può degenerare in una retta.}. % lang-it
Definiamo la curvatura con segno come $\pm 1 / r$, dove il segno $-$ corrisponde a una tangenza esterna. % lang-it
Allora% % lang-it
\begin{equation}
    \left(\pm \frac{1}{r_1} \pm  \frac{1}{r_2} \pm  \frac{1}{r_3} \pm  \frac{1}{r_4}\right)^2 = 2 \left(\frac{1}{r_1^2} + \frac{1}{r_2^2} + \frac{1}{r_3^2} + \frac{1}{r_4^2}\right),
\end{equation}
gdzie $r_4$ oznacza promień czwartego okręgu.% % lang-pl
dove $r_4$ indica il raggio del quarto cerchio.% % lang-it
\end{theorem}

Jeżeli jeden z okręgów (dajmy na to ten o promieniu $r_1$) staje się prostą, to jego zerowa krzywizna znika z równania. % lang-pl
Jeżeli dwa okręgi stają się dwiema prostymi, to pozostałem dwa muszą być do siebie przystające, co należy uznać za zdegenerowany przypadek.% % lang-pl
Se uno dei cerchi (ad esempio quello di raggio $r_1$) diventa una retta, la sua curvatura nulla scompare dall'equazione. % lang-it
Se due cerchi diventano due rette, gli altri due devono essere congruenti, il che va considerato un caso degenere.% % lang-it

Twierdzenie Kartezjusza odkryją na nowo Yamaji Nuchizumi (1751), Jakob Steiner (\emph{Fortsetzung der geometrischen Betrachtungen} z 1826), Philip Beecroft (\emph{Properties of circles in mutual contact} z 1842) i~wreszcie Frederick Soddy (\emph{The Kiss Precise} z 1936 jak opisane u Coxetera \cite[s. 29-31]{coxeter_1967}).% % lang-pl
Il teorema di Cartesio fu riscoperto da Yamaji Nuchizumi (1751), Jakob Steiner (\emph{Fortsetzung der geometrischen Betrachtungen}, 1826), Philip Beecroft (\emph{Properties of circles in mutual contact}, 1842) e infine Frederick Soddy (\emph{The Kiss Precise}, 1936, come descritto in Coxeter \cite[s. 29-31]{coxeter_1967}).% % lang-it
\index[persons]{Nuchizumi, Yamaji}%
\index[persons]{Steiner, Jakob}%
\index[persons]{Beecroft, Philip}%
\index[persons]{Soddy, Frederick}%

Ten ostatni poda uogólnienie do sfer, a rok później Thorold Gosset przeskoczy do przestrzeni dowolnego wymiaru ($\mathbb R^n$), z $n$ w miejsce czynnika $2$.% % lang-pl
Quest'ultimo fornì una generalizzazione alle sfere, e un anno dopo Thorold Gosset passò allo spazio di dimensione arbitraria ($\mathbb R^n$), con $n$ al posto del fattore $2$.% % lang-it
\index[persons]{Gosset, Thorold}%

Dowody korzystają z inwersji, wzoru Herona, wyznacznika Cayleya-Mengera, albo łańcuchów Pappusa i twierdzenia Vivianiego. % lang-pl
Patrz też $\Delta_{10}^6$.% % lang-pl
Le dimostrazioni utilizzano inversioni, la formula di Erone, il determinante di Cayley--Menger, oppure le catene di Pappo e il teorema di Viviani.% % lang-it
Si veda anche $\Delta_{10}^6$. % lang-it

\begin{example}
Dane są trzy styczne zewnętrznie okręgi o promieniach $r_1 = r_2 = r_3 = \sqrt{3}$. % lang-pl
Czwarty okrąg styczny do nich wszystkich ma wtedy krzywiznę $\sqrt 3 \pm 2$.% % lang-pl
Sono dati tre cerchi tangenti esternamente con raggi $r_1 = r_2 = r_3 = \sqrt{3}$. % lang-it
Il quarto cerchio tangente a tutti ha allora curvatura $\sqrt 3 \pm 2$.% % lang-it
\end{example}

Jeszcze nie skończyliśmy z Soddym. % lang-pl
W wierzchołkach trójkąta o obwodzie $2p = a + b + c$ wstawiamy okręgi o promieniach $p - a$, $p - b$, $p - c$. % lang-pl
Wtedy czwarty okrąg (zwany okręgiem Soddy'ego) z twierdzenia Kartezjusza ma promień% % lang-pl
Non abbiamo ancora finito con Soddy. % lang-it
Nei vertici di un triangolo di perimetro $2p = a + b + c$ inseriamo cerchi di raggi $p - a$, $p - b$, $p - c$. % lang-it
Allora il quarto cerchio (detto cerchio di Soddy) dal teorema di Cartesio ha raggio% % lang-it
\begin{equation}
    \frac{4R + r \pm 2p}{S},
\end{equation}
gdzie $R, r$ to promienie okręgu opisanego i wpisanego, zaś $S$ to pole powierzchni. % lang-pl
Napisze o tym Coxeter \cite[s. 29-31]{coxeter_1967}.% % lang-pl
dove $R, r$ sono i raggi del cerchio circoscritto e inscritto, e $S$ è l'area. % lang-it
Ne scrive Coxeter \cite[s. 29-31]{coxeter_1967}.% % lang-it

\begin{definition}[prosta Soddy'ego] % lang-pl
    Prostą przechodzącą przez środki okręgów Soddy'ego nazywamy prostą Soddy'ego. % lang-pl
\begin{definition}[retta di Soddy] % lang-it
    
    La retta che passa per i centri dei cerchi di Soddy si chiama retta di Soddy. % lang-it
\end{definition}

\begin{proposition}
Prosta Soddy'ego przechodzi przez środek okręgu wpisanego, punkt Gergonne'a i punkt Eppsteina (punkt, gdzie przecinają się trzy proste przechodzące przez trzy pary spośród sześciu punktów styczności czterech okręgów...)% % lang-pl
\index{punkt!Eppsteina}% % lang-pl
La retta di Soddy passa per il centro del cerchio inscritto, il punto di Gergonne e il punto di Eppstein (il punto di intersezione di tre rette che passano per tre coppie dei sei punti di tangenza dei quattro cerchi...).% % lang-it
\index{punto!di Eppstein}% % lang-it
\end{proposition}

Prosta Soddy'ego ma pewien związek z punktem de Lonchampsa, którego nie opiszemy.% % lang-pl
\index{punkt!de Longchampsa}% % lang-pl
La retta di Soddy ha una certa relazione con il punto di de Longchamps, che non descriveremo.% % lang-it
\index{punto!di de Longchamps}% % lang-it

%
%

\subsection{Problem Cramera-Castillona} % lang-pl
\subsection{Problema di Cramer-Castillon} % lang-it

Poniższy problem będzie mieć ciekawą historię. % lang-pl
Il seguente problema avrà una storia interessante. % lang-it

\index{zadanie!Cramera-Castillona}
\index{problema!di Cramer-Castillon}
\begin{problem}
\label{problem_cramera_castillona}%
    Dane jest $n$ punktów oraz okrąg. % lang-pl
    Znaleźć taki wielokąt wpisany w~okrąg, że na każdym jego boku (lub jego przedłużeniu) znajduje się dokładnie jeden z~danych punktów. % lang-pl
    %
    Sono dati $n$ punti e un cerchio. % lang-it
    Trovare un poligono inscritto nel cerchio tale che su ciascun lato (o sul suo prolungamento) si trovi esattamente uno dei punti assegnati. % lang-it
\end{problem}
% TODO: https://ems.press/content/serial-article-files/45288 ładny obrazek, związek z Urquhartem

Szczególny przypadek $n = 3$ współliniowych punktów zainteresuje Pappusa. % lang-pl
Il caso particolare di $n = 3$ punti allineati interesserà Pappo. % lang-it
\index[persons]{Pappus}% % lang-pl
\index[persons]{Pappo}% % lang-it
Pewien nieznany, ale stary geometra przedstawi problem dla dowolnych trzech punktów Gabrielowi Cramerowi, który w~1742 przekaże go dalej do Giovanniego Salveminiego\footnote{Giovanni Francesco Mauro Melchiorre Salvemini di Castiglione przyjmie w pewnym roku nazwisko od swojego miejsca urodzenia, Castiglione del Valdarno w Toskanii.}. % lang-pl
Un geometra ignoto, ma antico, presenterà il problema per tre punti arbitrari a Gabriel Cramer, che nel 1742 lo trasmetterà a Giovanni Salvemini\footnote{Giovanni Francesco Mauro Melchiorre Salvemini di Castiglione assunse a un certo punto il cognome dal suo luogo di nascita, Castiglione del Valdarno in Toscana.}. % lang-it
\index[persons]{Castiglione, Giovanni}%
\index[persons]{Cramer, Gabriel}%
Ten znajdzie geometryczne rozwiązanie już po śmierci Cramera, w~1776 roku; wyczyn powtórzą Euler (1783 rok), Lagrange, Malfatti, Lhuilier, Servois, Poncelet % lang-pl
Egli troverà una soluzione geometrica dopo la morte di Cramer, nel 1776; il risultato sarà poi ottenuto nuovamente da Eulero (1783), Lagrange, Malfatti, L'Huillier, Servois e Poncelet % lang-it
\index[persons]{Lagrange, Joseph Louis}%
\index[persons]{Malfatti, Gian Francesco}%
\index[persons]{L'Huillier, Simon Antoine Jean}%
\index[persons]{Servois, François-Joseph}%
\index[persons]{Poncelet, Jean-Victor}%
% https://cms.math.ca/wp-content/uploads/crux-pdfs/Crux_v9n04_Apr.pdf Crux Mathematicorum, Vol 9, No 4, page 125
i szesnastoletni Annibale\footnote{Carnot uzna, że Ottajano (miejsce urodzenia Giordana) jest tytułem szlacheckim, a~nie nazwą miejscowości; w~swoich publikacjach nazwie młodego matematyka właśnie Ottajano. Niestety, ale później będzie tak określany także w~kolejnych pracach naukowych.} Giordano. % lang-pl
e dal sedicenne Annibale\footnote{Carnot ritenne che Ottajano (luogo di nascita di Giordano) fosse un titolo nobiliare anziché il nome di una località; nelle sue pubblicazioni chiamò quindi il giovane matematico Ottajano. Purtroppo questa denominazione verrà ripresa anche in lavori successivi.} Giordano. % lang-it
\index[persons]{Giordano, Annibale}%
Carnot uprości analityczne rozwiązanie Lagrange'a i~uogólni je do $n \ge 3$ punktów. % lang-pl
Carnot semplificherà la soluzione analitica di Lagrange e la generalizzerà al caso di $n \ge 3$ punti. % lang-it
\index[persons]{Carnot, Lazare}%

% TODO: https://en.wikipedia.org/wiki/Cramer-Castillon_problem

%
%

\index{zadanie!Monge'a}
\pol{\subsection{Problem Monge'a}}
\ita{\subsection{Problema di Monge}}

\begin{problem}[\pol{Monge'a?}\ita{di Monge?}]
    \pol{Okrąg, który przecina trzy okręgi $\Gamma_1$, $\Gamma_2$, $\Gamma_3$ pod kątem prostym.}
    \ita{Costruire un cerchio che intersechi ortogonalmente tre cerchi $\Gamma_1$, $\Gamma_2$, $\Gamma_3$.}
\end{problem}

% https://mathworld.wolfram.com/MongesProblem.html
\pol{Dla każdej pary okręgów $\Gamma_i$, $\Gamma_j$ znajdujemy oś potęgową; jeśli trzy osie przecinają się w punkcie $O$ leżącym na zewnątrz okręgów $\Gamma_i$, to jest to środek szukanego okręgu.}
\ita{Per ogni coppia di cerchi $\Gamma_i$, $\Gamma_j$ si determina l'asse radicale; se i tre assi si intersecano in un punto $O$ situato all'esterno dei cerchi $\Gamma_i$, allora esso è il centro del cerchio cercato.}
\pol{Promieniem szukanego okręgu jest odcinek styczny do $\Gamma_i$ oraz przechodzący przez $O$.}
\ita{Il raggio del cerchio cercato è dato dal segmento tangente a $\Gamma_i$ passante per $O$.}
\pol{Jeśli jednak środek okręgu leży wewnątrz któregoś okręgu albo osie nie przecinają się, to problem nie ma rozwiąania.}
\ita{Se invece il centro del cerchio giace all'interno di uno dei cerchi oppure gli assi non si intersecano, allora il problema non ha soluzione.}

%
%

\subsection{Zadanie Malfattiego}
\index{zadanie!Malfattiego}

W 1803 roku Malfatti \cite{malfatti_1803} zainspirowany pewnym praktycznym zagadnieniem (wycinanie walców z graniastosłupa) postawi następujący problem:
\index[persons]{Malfatti, Gian Francesco}%

\begin{problem}[zadanie Malfattiego]
	\label{malfatti_problem}
	\index{zadanie!Malfattiego}%
	Dany jest trójkąt $\triangle ABC$.
	Skonstruować takie trzy parami styczne okręgi $\Gamma_A, \Gamma_B, \Gamma_C$, że okrąg $\Gamma_A$ (odpowiednio: $\Gamma_B$, $\Gamma_C$) jest wpisany w~kąt $\angle A$ (odpowiednio: $\angle B$, $\angle C$).
\end{problem}

% https://www.desmos.com/calculator/mqzextwkad?lang=pl
\begin{figure}[H] \centering
\begin{comment}
\begin{tikzpicture}[scale=.5]
\tkzDefPoints{0/0/A,10/2/B,6/7/C}
\tkzDefPoints{4.43012726/2.59439459/Oa}
\tkzDefCircle[R](Oa,1.67519375895) \tkzGetPoint{Oaa}
\tkzDrawCircle[line width=0.2mm](Oa,Oaa)

\tkzDefPoints{7.48168986/2.91734309/Ob}
\tkzDefCircle[R](Ob,1.39341015784) \tkzGetPoint{Obb}
\tkzDrawCircle[line width=0.2mm](Ob,Obb)

\tkzDefPoints{5.96721113/5.06490116/Oc}
\tkzDefCircle[R](Oc,1.23445046858) \tkzGetPoint{Occ}
\tkzDrawCircle[line width=0.2mm](Oc,Occ)

\tkzLabelPoint(A){$A$}
\tkzLabelPoint[anchor=center](Oa){$\Gamma_A$}
\tkzLabelPoint(B){$B$}
\tkzLabelPoint[anchor=center](Ob){$\Gamma_B$}
\tkzLabelPoint[above](C){$C$}
\tkzLabelPoint[anchor=center](Oc){$\Gamma_C$}
\tkzDrawPolygon[line width=0.3mm](A,B,C)
\end{tikzpicture}
\end{comment}
\caption{Trzy okręgi Malfattiego}
\end{figure}

Problem będzie rozważany na długo przed Malfattim, zajmie się nim Ajima Naonobu\footnote{Matematyk japoński, przypisze się mu wprowadzenie rachunku różniczkowo-całkowego do matematyki japońskiej.} w~XVIII wieku, a~jeszcze wcześniej Gilio de Cecco da Montepulciano w~rękopisie z~1384 roku.
\index[persons]{Ajima, Naonobu}%
\index[persons]{de Cecco da Montepulciano, Gilio}%

Malfatti wyprowadzi co następuje.
Niech $p$ będzie połową obwodu trójkąta, $r$ będzie promieniem okręgu wpisanego w~ten trójkąt zaś $d_A$, $d_B$, $d_C$ odległościami wierzchołków $A, B, C$ od środka tego okręgu.
Wtedy promienie okręgów Malfattiego wyrażają się wzorami
\begin{align}
	r_A & = \frac r 2 \cdot {\frac {s-r+d_A-d_B-d_C}{p-a}}, \\
	r_B & = \frac r 2 \cdot {\frac {s-r+d_B-d_A-d_C}{p-b}}, \\
	r_C & = \frac r 2 \cdot {\frac {s-r+d_C-d_A-d_B}{p-c}}.
\end{align}

Prostą konstrukcję okręgów opartą na dwustycznych zawdzięczymy Steinerowi \cite{steiner_1826} w~1826 roku;
\index[persons]{Steiner, Jakob}%
inne rozwiązania podadzą Lehmus \cite{lehmus_1819}, Catalan \cite{catalan_1846}, Adams \cite{adams_1846}, Derousseau \cite{derousseau_1895}, Pampuch \cite{pampuch_1904}.
% TODO: po poprawie bibliografii, podać tu index persons

(O~problemie napiszą też Bogdańska, Neugebauer \cite[s. 102]{neugebauer_2018}).

Malfatti postawi tak naprawdę inny problem: znalezienia trzech rozłącznych kół zawartych w~trójkącie, których suma pól jest maksymalna i~błędnie założy, że opisane wyżej okręgi stanowią rozwiązanie.
Pomyłkę zauważą najpierw bez dowodu Lob, Richmond \cite{lob_richmond_1930} w~1930 roku: z trójkąta równobocznego można wyciąć zachłannie kolejno trzy koła, ich łączna powierzchnia jest większa od powierzchni kół znalezionych przez Malfattiego o 1\%.
\index[persons]{Richmond, ?}%
\index[persons]{Lob, ?}%
Howard Eves powtórzy to dla stromych trójkątów równoramiennych o bardzo wąskiej podstawie i dużej wysokości około 1946 roku.
\index[persons]{Eves, Howard}%
% https://en.wikipedia.org/w/index.php?title=Howard_Eves&diff=831382284&oldid=750910758
Goldberg \cite{goldberg_1967} wykaże, że domniemanie Malfattiego nie daje nigdy kół o maksymalnej łącznej powierzchni.
Ostatnie słowo należy zaś do Zalgallera, Losa \cite{zalgaller_los_1992}, którzy znajdą trzy koła rozwiązujące problem Malfattiego w dowolnym trójkącie.
% TODO: Goldberg M., On the original Malfatti problem, Math. Mag. 40 (1967), 241-247.
\index[persons]{Zalgaller, VA?}%
\index[persons]{Los, GA?}%
% TODO: Zalgaller V.A., Los’ G.A., Solution of the Malfatti problem, Ukrain. Geom. Sb. 35 (1992), 14-33 (ang. J. Math. Sci. 72 (1994), 3163-3177).
% TODO: po poprawie bibliografii, podać tu index persons
% TODO: Lob, H.; Richmond, H. W. (1930), "On the Solutions of Malfatti's Problem for a Triangle", Proceedings of the London Mathematical Society, 2nd ser., 30 (1): 287-304, doi:10.1112/plms/s2-30.1.287.

Kryształowa kula nie potrafi przewidzieć, kto oceni, czy algorytm zachłanny zawsze znajduje $n \ge 4$ rozłącznych kół w trójkącie o maksymalnej łącznej powierzchni.

(O więcej niż jednym okręgu wpisanym w trójkąt pisaliśmy w podpodsekcji \ref{sssection_6_7_9_circles}).


\subsection{Punkty Crelle'a-Brocarda} % lang-pl
\subsection{Punti di Crelle-Brocard} % lang-it
\index{punkt!Crelle'a-Brocarda}% % lang-pl
\index{punto!di Crelle-Brocard}% % lang-it

Punkty Crelle'a-Brocarda (\ref{punkty_brocarda}) doczekają się eleganckiej konstrukcji cyrklem i linijką. % lang-pl
I punti di Crelle-Brocard (\ref{punkty_brocarda}) avranno un'elegante costruzione con riga e compasso. % lang-it
Aby otrzymać pierwszy punkt, poprowadzimy okrąg przez wierzchołki $A$ i $B$, styczny do boku $BC$ w~$B$; symetrycznie okrąg przez $B$ i $C$, styczny do boku $CA$ w~$C$; oraz okrąg przez $C$ i $A$, styczny do boku $AB$ w~$A$. % lang-pl
Per ottenere il primo punto, tracceremo una circonferenza per i vertici $A$ e $B$, tangente al lato $BC$ in $B$; simmetricamente una circonferenza per $B$ e $C$, tangente al lato $CA$ in $C$; e una circonferenza per $C$ e $A$, tangente al lato $AB$ in $A$. % lang-it
Te trzy okręgi przetną się w jednym punkcie -- właśnie pierwszym punkcie Crelle'a-Brocarda. % lang-pl
Queste tre circonferenze si incontreranno in un unico punto -- proprio il primo punto di Crelle-Brocard. % lang-it
Drugi punkt otrzymamy analogicznie, tylko odwracając kierunek styczności: okrąg przez $A$ i $B$ styczny do $AC$ w~$A$, okrąg przez $B$ i $C$ styczny do $AB$ w~$B$ oraz okrąg przez $C$ i $A$ styczny do $BC$ w~$C$. % lang-pl
Il secondo punto si ottiene in modo analogo, invertendo semplicemente il verso della tangenza: una circonferenza per $A$ e $B$ tangente ad $AC$ in $A$, una circonferenza per $B$ e $C$ tangente ad $AB$ in $B$ e una circonferenza per $C$ e $A$ tangente a $BC$ in $C$. % lang-it
% https://en.wikipedia.org/wiki/Brocard_points#Construction

\subsection{Lista Wernicksa} % lang-pl
\subsection{Lista di Wernick} % lang-it

\emph{Sixteen key points of a triangle are its vertices, the midpoints of its sides, the feet of its altitudes, the feet of its internal angle bisectors, and its circumcenter, centroid, orthocenter, and incenter. These can be taken three at a time to yield 139 distinct nontrivial problems of constructing a triangle from three points.[12] Of these problems, three involve a point that can be uniquely constructed from the other two points; 23 can be non-uniquely constructed (in fact for infinitely many solutions) but only if the locations of the points obey certain constraints; in 74 the problem is constructible in the general case; and in 39 the required triangle exists but is not constructible.} % lang-pl
\emph{Sedici punti notevoli di un triangolo sono i suoi vertici, i punti medi dei lati, i piedi delle altezze, i piedi delle bisettrici interne e inoltre il circocentro, il baricentro, l'ortocentro e l'incentro. Considerandoli a gruppi di tre si ottengono 139 distinti problemi non banali di costruzione di un triangolo a partire da tre punti.[12] Di questi problemi, tre coinvolgono un punto che può essere costruito univocamente a partire dagli altri due; 23 ammettono costruzioni non univoche (in effetti con infinite soluzioni), ma solo se le posizioni dei punti soddisfano determinati vincoli; in 74 casi il problema è costruibile nel caso generale; mentre in 39 casi il triangolo richiesto esiste ma non è costruibile.} % lang-it
%  => https://en.wikipedia.org/wiki/Straightedge_and_compass_construction
% => Zetel s. 32 rtrójkąt ze środkowych
% todo: zadanie 1.15 spodki dwusiecznych gdzeiś wcisnąć?

% TODO: https://en.wikipedia.org/wiki/Malfatti_circles

%
%

\section{Wielokąty foremne} % lang-pl
\section{Poligoni regolari} % lang-it
\index{wielokąt!foremny}

\begin{problem}
    Skonstrować trójkąt równoboczny wpisany w okrąg, którego środek nie jest znany. \hfill \emph{(7 kroków)} % lang-pl
    Costruire un triangolo equilatero inscritto in un cerchio di cui non è noto il centro. \hfill \emph{(7 passi)} % lang-it
\end{problem}

\begin{problem}
    Skonstrować kwadrat. \hfill \emph{(9 kroków)} % lang-pl
    Costruire un quadrato. \hfill \emph{(9 passi)} % lang-it
\end{problem}

Konstrukcja trójkąta równobocznego oraz kwadratu będzie znana dobrze w~starożytności i~nie wiadomo, kto pierwszy na nią wpadnie. % lang-pl
La costruzione del triangolo equilatero e del quadrato sarà ben nota nell'antichità e non si sa chi l'abbia scoperta per primo. % lang-it
Pojawią się u Euklidesa (I.1), (I.46). % lang-pl
Esse compaiono negli \emph{Elementi} di Euclide (I.1), (I.46). % lang-it

\begin{problem}
    Skonstrować pięciokąt foremny. % lang-pl
    Costruire un pentagono regolare. % lang-it
\end{problem}

Pięciokąt foremny jest pierwszym wielokątem foremnym, którego konstrukcja nie jest oczywista. % lang-pl
Il pentagono regolare è il primo poligono regolare la cui costruzione non è immediata. % lang-it
Różne pomysły podadzą Euklides (IV.11), Ptolemeusz w~Almageście (łatwiej będzie jednak przeczytać Guzickiego \cite[s. 286-287]{guzicki_2021}). % lang-pl
Diverse costruzioni saranno proposte da Euclide (IV.11) e da Tolomeo nell'\emph{Almagesto} (anche se sarà più semplice leggere Guzicki \cite[p. 286-287]{guzicki_2021}). % lang-it
\index[persons]{Ptolemeusz}% % lang-pl
\index[persons]{Tolomeo}% % lang-it
\index{Almagest}%
Nie pominie jej Hartshorne \cite[s. 45-49]{hartshorne2000} w swoim ,,przedłużeniu'' Elementów z komentarzem, że dowód Euklidesa poprawności konstrukcji tworzy standard pięknego dowodu matematycznego. % lang-pl
Hartshorne \cite[p. 45-49]{hartshorne2000} la include nella sua sorta di ``continuazione'' degli \emph{Elementi}, osservando che la dimostrazione di Euclide costituisce un modello di elegante dimostrazione matematica. % lang-it
Jeszcze inne eleganckie rozwiązanie (opisane w~$\Delta_{24}^{10}$) znajdzie Yosifusa Hirano, japoński miłośnik matematyki z XIX wieku, a opisze jego przyjaciel Kawakita Chorin w~manuskrypcie ,,\emph{Sanpo Jyojutu Kaigi}'' (rozwiązania Sanpo Jyojutu). % lang-pl
Un'altra elegante soluzione (descritta in $\Delta_{24}^{10}$) sarà trovata da Yosifusa Hirano, appassionato giapponese di matematica del XIX secolo, e verrà descritta dal suo amico Kawakita Chorin nel manoscritto \emph{Sanpo Jyojutu Kaigi} (Soluzioni del Sanpo Jyojutu). % lang-it
\index[persons]{Hirano, Yosifusa}%
Konstrukcja Hirano to kolejny powód, by sięgnąć po książkę Fukagawy, Pedoe \cite{fukagawa_1989}. % lang-pl
La costruzione di Hirano è un'ulteriore ragione per consultare il libro di Fukagawa e Pedoe \cite{fukagawa_1989}. % lang-it

Jeśli mamy zadany jeden z jego boków, konstrukcja wymaga 11 kroków. % lang-pl
Se è assegnato uno dei lati, la costruzione richiede 11 passi. % lang-it
% Hartshorne s. 51

\begin{problem}
    Skonstruować wielokąt foremny o $n$ kątach. % lang-pl
    Costruire un poligono regolare con $n$ lati. % lang-it
\end{problem}

Poza wyżej wymienionymi, Euklides potrafi skonstruować sześciokąt foremny (IV.15) albo, po wykreśleniu trójkąta i pięciokąta opisanych wewnątrz tego samego okręgu, piętnastokąt (IV.16). % lang-pl
Oltre ai casi già menzionati, Euclide sa costruire un esagono regolare (IV.15) e, dopo aver costruito nello stesso cerchio un triangolo e un pentagono, anche un pentadecagono (IV.16). % lang-it
\index[persons]{Euklides}% % lang-pl
\index[persons]{Euclide}% % lang-it
Mając wielokąt o $n$ bokach, bisekcja kąta środkowego pozwala podwoić liczbę boków do $2n$. % lang-pl
Dato un poligono di $n$ lati, la bisezione dell'angolo al centro permette di raddoppiare il numero dei lati fino a $2n$. % lang-it
Możemy więc uznać, że Euklides wykreśli wielokąty, które mają % lang-pl
Possiamo quindi considerare che Euclide sia in grado di costruire poligoni con % lang-it
\begin{equation}
    3, 4, 5, 6, 8, 10, 12, 15, 16, 20, \ldots
\end{equation}
boków. % lang-pl
lati. % lang-it
Lista zostanie rozszerzona dopiero w 1796 roku przez Gaussa: pokaże, że siedemnastokąt foremny jest konstruowalny. % lang-pl
L'elenco verrà ampliato soltanto nel 1796 da Gauss, che mostrerà che il poligono regolare di 17 lati è costruibile. % lang-it
\index{Gauss, Carl Friedrich}%
\index{wielokąt!siedemnastokąt}% % lang-pl
\index{poligono!regolare di 17 lati}% % lang-it
Będzie tak zadowolony z tego wyniku, że zażyczy sobie wyryć właśnie tę figurę na swoim grobie. % lang-pl
Sarà così soddisfatto di questo risultato da desiderare che tale figura venga incisa sulla sua tomba. % lang-it
Faktycznej konstrukcji dokona Johannes Erchinger w 1800 roku. % lang-pl
La costruzione effettiva sarà realizzata da Johannes Erchinger nel 1800. % lang-it
\index[persons]{Erchinger, Johannes}%
Hartstone \cite[s. 250-259]{hartshorne2000} poświęca tej konstrukcji całą sekcję; Guzicki nie szczędzi rachunków \cite[s. 290-299]{guzicki_2021}. % lang-pl
Hartshorne \cite[p. 250-259]{hartshorne2000} dedica a questa costruzione un'intera sezione; Guzicki non risparmia i dettagli di calcolo \cite[p. 290-299]{guzicki_2021}. % lang-it
Siedemnastokąt pojawi się też u Evesa \cite[s. 185-191]{eves1_1972} i Coxetera \cite[s. 42, 43]{coxeter_1967}. % lang-pl
L'eptadecagono compare anche in Eves \cite[p. 185-191]{eves1_1972} e in Coxeter \cite[p. 42, 43]{coxeter_1967}. % lang-it

W ogólności, mamy: % lang-pl
Più in generale, vale il seguente risultato: % lang-it
\begin{theorem}
[Gaussa-Wantzla] % lang-pl
[di Gauss-Wantzel] % lang-it
    \label{gauss_wantzel}%
    Wielokąt foremny o $n$ bokach jest konstruowalny przy pomocy cyrkla i~linijki bez podziałki wtedy i tylko wtedy, kiedy $n$ jest postaci % lang-pl
    Un poligono regolare di $n$ lati è costruibile con riga non graduata e compasso se e solo se $n$ è della forma % lang-it
    \begin{equation}
        n = 2^r p_1 \cdot \ldots \cdot p_s,
    \end{equation}
    $r, s \ge 0$, gdzie $p_i$ są różnymi liczbami pierwszymi postaci $2^{2^k} + 1$, na przykład: $3$, $5$, $17$, $257$, $65537$. % lang-pl
    $r, s \ge 0$, dove i $p_i$ sono distinti numeri primi della forma $2^{2^k}+1$, per esempio: $3$, $5$, $17$, $257$, $65537$. % lang-it
    \index{liczba pierwsza!Fermata}% % lang-pl
    \index{numero primo!Fermat}% % lang-it
\end{theorem}

Gauss znajdzie warunek wystarczający i napisze bez dowodu w~\emph{Disquisitiones Arithmeticae} (po polsku byłyby to Dociekaniach arytmetyczne, 1801), że jest też konieczny. % lang-pl
Gauss troverà la condizione sufficiente e scriverà senza dimostrazione nelle \emph{Disquisitiones Arithmeticae} (1801) che essa è anche necessaria. % lang-it
Brakujący dowód doda Pierre Laurent Wantzel w 1837 roku. % lang-pl
La dimostrazione mancante sarà fornita da Pierre Laurent Wantzel nel 1837. % lang-it
\index[persons]{Wantzel, Pierre Laurent}%
Spisany współczesnym językiem znajdzie się u Hartstone'a \cite[s. 258]{hartshorne2000}; wzmiankować go będzie Guzicki \cite[s. 285]{guzicki_2021}. % lang-pl
Una presentazione in linguaggio moderno si trova in Hartshorne \cite[p. 258]{hartshorne2000}; Guzicki ne fa menzione in \cite[p. 285]{guzicki_2021}. % lang-it

Twierdzenie Gaussa zachęca do znalezienia przepisu na wielokąty foremne o większej liczbie boków. % lang-pl
Il teorema di Gauss incoraggia a cercare costruzioni di poligoni regolari con un numero ancora maggiore di lati. % lang-it
W 1822 roku Magnus Georg von Paucker\footnote{Właściwie Магнус-Георг Андреевич Паукер, urodzon w Simunie, wtedy części Imperium Rosyjskiego, później Łotwy.} \cite{paucker_1822}, a po dziesięciu latach także Friederich Julius Richelot \cite{richelot_1832} znajdą konstrukcję $257$-kąta foremnego. % lang-pl
Nel 1822 Magnus Georg von Paucker\footnote{In realtà Магнус-Георг Андреевич Паукер, nato a Simuna, allora parte dell'Impero Russo e oggi in Estonia.} \cite{paucker_1822}, e dieci anni più tardi anche Friederich Julius Richelot \cite{richelot_1832}, troveranno una costruzione del poligono regolare di $257$ lati. % lang-it
\index[persons]{Paucker@von Paucker, Magnus Georg}%
\index[persons]{Richelot, Friederich Julius}%
\index[persons]{Hermes, Johann(es?)}%
Ich cierpliwość przyćmiewa konstrukcja $65\,537$-kąta Johanna Hermesa z 1896 roku, na którą poświęci dziesięć lat swojego życia (i dwieście stron rękopisu). % lang-pl
La loro pazienza sarà però superata dalla costruzione del poligono regolare di $65\,537$ lati realizzata da Johann Hermes nel 1896, alla quale dedicò dieci anni della sua vita e duecento pagine di manoscritto. % lang-it

W 2025 roku nie będzie jeszcze wiadomo, ile jest liczb pierwszych Fermata, więc ludzkość będzie w stanie wykreślić wielokąty o nieparzyście wielu bokach dokładnie wtedy, gdy liczba tych boków będzie dzielnikiem $2^{32} - 1 = 4\,294\,967\,295$. % lang-pl
Nel 2025 non si saprà ancora quante siano le prime di Fermat, per cui l'umanità sarà in grado di costruire poligoni con un numero dispari di lati esattamente quando tale numero divide $2^{32}-1=4\,294\,967\,295$. % lang-it
Dla oszczędności papieru i tuszu darujemy sobie wypisywanie tych dzielników. % lang-pl
Per risparmiare carta e inchiostro eviteremo di elencare tutti questi divisori. % lang-it

% TODO: \todofoot{Coxeter, s. 26}

\subsection{Konstrukcje przybliżone} % lang-pl
\subsection{Costruzioni approssimate} % lang-it

Dwie konstrukcje przybliżone pięciokąta foremnego podadzą Leonardo da Vinci (gdzie kąty środkowe mają miarę $71.6350$ oraz $72.5475$ zamiast dokładnie $72$ stopni) oraz Dürer; obydwie można znaleźć u Guzickiego \cite[s. 287-290]{guzicki_2021}. % lang-pl
Due costruzioni approssimate del pentagono regolare saranno proposte da Leonardo da Vinci (nelle quali gli angoli al centro misurano rispettivamente $71.6350$ e $72.5475$ gradi invece di $72$) e da Dürer; entrambe si trovano in Guzicki \cite[p. 287-290]{guzicki_2021}. % lang-it

\subsection{Konstrukcje z niestandardowymi przyborami} % lang-pl
\subsection{Costruzioni con strumenti non standard} % lang-it

Jeżeli mamy do dyspozycji trysektor kąta, potrafimy zbudować więcej wielokątów foremnych niż sugeruje to twierdzenie \ref{gauss_wantzel}: % lang-pl
Se disponiamo di un trisettore d'angolo, possiamo costruire più poligoni regolari di quanti ne consenta il teorema \ref{gauss_wantzel}: % lang-it
\index{trysekcja!kąta}% % lang-pl
\index{trysektor kąta}% % lang-pl
\index{trisettore d'angolo}% % lang-it
\index{???}% % lang-it

\begin{proposition}
    Wielokąt foremny o $n$ bokach jest konstruowalny przy pomocy cyrkla, linijki bez podziałki oraz trysektora kąta wtedy i tylko wtedy, kiedy $n$ jest postaci % lang-pl
    Un poligono regolare di $n$ lati è costruibile con compasso, riga non graduata e trisettore d'angolo se e solo se $n$ è della forma % lang-it
    \begin{equation}
        n = 2^r 3^q p_1 \cdot \ldots \cdot p_s,
    \end{equation}
    $q, r, s \ge 0$, gdzie $p_i$ są różnymi liczbami pierwszymi Pierponta, postaci $2^t3^u + 1$. % lang-pl
    $q, r, s \ge 0$, dove i $p_i$ sono distinti numeri primi di Pierpont della forma $2^t3^u + 1$. % lang-it
    \index{liczba pierwsza!Pierponta}% % lang-pl
    \index{numero primo!Pierpont}% % lang-it
\end{proposition}

Dowód pokaże Hartshorne \cite[s. 276]{hartshorne2000}. % lang-pl
% TODO: włoski brakuje

To są dokładnie wielokąty, które można skonstruować używając stożkowych, albo równoważnie origami. % lang-pl
Questi sono esattamente i poligoni che si possono costruire usando le coniche oppure, in modo equivalente, l'origami. % lang-it
Lista nowych nieparzystych liczb $n$, który dostarcza ten fakt, to \textbf{7}, $9$, \textbf{13}, \textbf{19}, $21$, $27$, $35$, \textbf{37}, $39$, $45$, $57$, $63$, $65$, \textbf{73}, $81$, $91$, $95$, \textbf{97}, $105$, \textbf{109}, $111$, $117$, $119$, $133$, $135$, $153$, \textbf{163}, $171$, $185$, $189$, \textbf{193}, $195$, $219$, $221$, $243$, $247$, \textbf{257}, $259$, $273$, $285$, $291\ldots$ (pogrubiono liczby pierwsze). % lang-pl
L'elenco dei nuovi valori dispari di $n$ forniti da questo risultato è \textbf{7}, $9$, \textbf{13}, \textbf{19}, $21$, $27$, $35$, \textbf{37}, $39$, $45$, $57$, $63$, $65$, \textbf{73}, $81$, $91$, $95$, \textbf{97}, $105$, \textbf{109}, $111$, $117$, $119$, $133$, $135$, $153$, \textbf{163}, $171$, $185$, $189$, \textbf{193}, $195$, $219$, $221$, $243$, $247$, \textbf{257}, $259$, $273$, $285$, $291\ldots$ (i numeri primi sono evidenziati in grassetto). % lang-it
% gdzie 25???

Szczegóły zapewnia artykuł Gleasona \cite{gleason_1988}. % lang-pl
I dettagli sono forniti nell'articolo di Gleason \cite{gleason_1988}. % lang-it
\index[persons]{Gleason, Andrew Mattei}%
Konstrukcję siedmiokąta przy użyciu znaczonej linijki opisuje Hartshorne \cite[s. 259-270]{hartshorne2000}. % lang-pl
La costruzione dell'eptagono mediante una riga graduata è descritta da Hartshorne \cite[p. 259-270]{hartshorne2000}. % lang-it
% TODO: https://www.deltami.edu.pl/2015/10/siedmiokata-foremnego-nie-mozna-skonstruowac-cyrklem-i-linijka/
\index{linijka!znaczona}% % lang-pl
\index{riga graduata}% % lang-it % TODO: auto translated

%
% eksperymentalne tłumaczenie na włoski

%

\section{Wyrocznia w Delfach} % lang-pl

\section{L'oracolo di Delfi} % lang-it
\index{wyrocznia w Delfach}

Przytoczymy teraz trzy problemy geometryczne, zaprzątające głowy przez blisko dwa tysiąclecia. % lang-pl

Wszystkie trzy znajdziemy też u Evesa \cite[s. 156]{eves1_1972}. % lang-pl
Presenteremo ora tre problemi geometrici che hanno occupato le menti dei matematici per quasi due millenni. % lang-it
Tutti e tre sono trattati anche da Eves \cite[p. 156]{eves1_1972}. % lang-it


Według legendy wyrocznia na wyspie Delos poradzi, by dla powstrzymania panującej tam zarazy trzeba dwukrotnie powiększyć ołtarz ofiarny, nie zmieniając jego sześciennego kształtu. % lang-pl
Secondo la leggenda, l'oracolo dell'isola di Delo consigliò, per fermare una pestilenza che vi infuriava, di raddoppiare il volume dell'altare sacrificale senza modificarne la forma cubica. % lang-it
\index{podwojenie sześcianu}

\begin{problem}[problem delijski] % lang-pl
\begin{problem}[problema di Delo] % lang-it
    \label{problem_delijski}%
    Zbudować sześcian o objętości dwa razy większej niż sześcian dany. % lang-pl
    
    Costruire un cubo avente volume doppio rispetto a un cubo assegnato. % lang-it
    \index{problem!delijski}
\end{problem}

https://www.deltami.edu.pl/2017/03/analiza-starozytnych-i-cyprian-norwid/

O problemie następne pokolenia dowiedzą się między innymi z prac Teona ze Smyrny. % lang-pl

Le generazioni successive conosceranno il problema anche grazie agli scritti di Teone di Smirne. % lang-it
\index[persons]{Teon ze Smyrny}% % lang-pl
\index[persons]{Teone di Smirne}% % lang-it
% https://en.wikipedia.org/wiki/Theon_of_Smyrna  It is also one of the sources of our knowledge of the origins of the classical problem of Doubling the cube.[5]
Choć odkryte zostanie wiele konstrukcji przybliżonych albo takich, które wykorzystują poza cyrklem i~linijką dodatkowe narzędzia (Hipokrates razem z Menaichmosem i hiperbola z parabolą, których nie da się wykreślić; Diokles i cysoida; Nikomedes i konchoida, Hippiasz z Elidy i kwadratrysa), to mimo upływu setek lat nie zostanie zrobiony żaden postęp. % lang-pl

Sebbene vengano scoperte molte costruzioni approssimate o costruzioni che utilizzano strumenti aggiuntivi oltre a riga e compasso (Ippocrate e Menecmo con l'iperbole e la parabola; Diocle con la cissoide; Nicomede con la concoide; Ippia di Elide con la quadratrice), nonostante il trascorrere dei secoli non si farà alcun progresso decisivo. % lang-it
\index[persons]{Hipokrates z Chios}% % lang-pl
\index[persons]{Menaichmos}% % lang-pl
\index{hiperbola}% % lang-pl
\index{parabola}%
\index[persons]{Diokles}% % lang-pl
\index{cysoida}% % lang-pl
\index[persons]{Nikomedes}% % lang-pl
\index{konchoida}% % lang-pl
\index[persons]{Hippiasz z Elidy}% % lang-pl
\index[persons]{kwadratrysa}% % lang-pl
\index[persons]{Ippocrate di Chio}% % lang-it
\index[persons]{Menecmo}% % lang-it
\index{iperbole}% % lang-it
\index[persons]{Diocle}% % lang-it
\index{cissoide}% % lang-it
\index[persons]{Nicomede}% % lang-it
\index{concoide}% % lang-it
\index[persons]{Ippia di Elide}% % lang-it
\index[persons]{quadratrice}% % lang-it

\begin{problem}[trysekcja kąta] % lang-pl
\begin{problem}[trisezione dell'angolo] % lang-it
    \label{trysekcja_kata}%
    Podzielić dowolny kąt na trzy równe części. % lang-pl
    
    Dividere un angolo arbitrario in tre parti uguali. % lang-it
    \index{trysekcja!kąta}
\end{problem}
% TODO: % TODO: Coxeter, s. 28
% TODO: https://en.wikipedia.org/wiki/Square_trisection
% https://en.wikipedia.org/wiki/Angle_trisection

https://www.deltami.edu.pl/2020/11/trysekcja-kata-w-geometrii-kartezjusza/

https://www.deltami.edu.pl/2015/09/zabawy-w-kacie/

Tu też pojawią się próby użycia różnorakich przyrządów. % lang-pl

Archmiedes zaproponuje dodanie podziałki do linijki, co okaże się wystarczające do wykreślenia kątów. % lang-pl
Anche qui verranno tentati approcci mediante strumenti di vario genere. % lang-it

Archimede proporrà di aggiungere una graduazione alla riga, il che si rivelerà sufficiente per la trisezione degli angoli. % lang-it
\index[persons]{Archimedes} % lang-pl
\index[persons]{Archimede} % lang-it
W średniowieczu zostanie zauważone, że zagadnienie trysekcji kąta prowadzi do rozwiązania równania algebraicznego stopnia trzeciego postaci % lang-pl

Nel Medioevo si osserverà che il problema della trisezione conduce alla risoluzione di un'equazione algebrica di terzo grado della forma % lang-it
\begin{equation}
    x^3 - 3 x - 2 \cos 3\alpha = 0.
\end{equation}
W szczególności, jeśli $\alpha = \pi / 9$, to pierwiastki nie są konstruowalne. % lang-pl

Patrz też do Coxetera \cite[s. 44]{coxeter_1967}. % lang-pl
In particolare, se $\alpha = \pi / 9$, le radici non sono costruibili. % lang-it
Si veda anche Coxeter \cite[p. 44]{coxeter_1967}. % lang-it
\begin{proposition} % lang-it
    La soluzione del problema di Delo e la trisezione dell'angolo non sono realizzabili con riga e compasso. % lang-it
\end{proposition} % lang-it


\begin{proposition} % lang-pl
    Rozwiązanie problemu delijskiego oraz trysekcja kąta są niewykonalne za pomocą cyrkla i linijki. % lang-pl
\end{proposition} % lang-pl
La prima dimostrazione di questo fatto fu data da Pierre Wantzel \cite{wantzel_1837} nel 1837. % lang-it

Pierwszy dowód tego faktu poda Pierre Wantzel \cite{wantzel_1837} w 1837 roku. % lang-pl
Być może to będzie powodem, dla którego dopiero w 1899 roku Frank Morley \cite{morley_1907} odkryje: % lang-pl
Forse anche per questo motivo soltanto nel 1899 Frank Morley \cite{morley_1907} scoprirà il seguente risultato: % lang-it
\index[persons]{Morley, Frank}

\begin{theorem}[Morleya]
    Punkty przeciecięcia tych spośród trójsiecznych kątów trójkąta, które sąsiadują z~którymś z~boków trójkąta, są wierzchołkami trójkąta równobocznego. % lang-pl
    
    I punti d'intersezione delle trisettrici degli angoli di un triangolo adiacenti ai suoi lati sono i vertici di un triangolo equilatero. % lang-it
    \index{trójkąt!Morleya}%
    \index{twierdzenie!Morleya}%
\end{theorem}

W tym sformułowaniu nie ma mowy, że chodzi o kąty wewnętrzne -- ponieważ twierdzenie jest prawdziwe także dla kątów zewnętrznych, z prawie takim samym dowodem; napisze o tym Kordos w $\Delta_{13}^{11}$. % lang-pl

Trójkąt równoboczny, jaki powstaje, ma bok długości % lang-pl
Nell'enunciato non si specifica che si tratta degli angoli interni, poiché il teorema è vero anche per gli angoli esterni, con una dimostrazione quasi identica; ne scrive Kordos in $\Delta_{13}^{11}$. % lang-it

Il triangolo equilatero ottenuto ha lato di lunghezza % lang-it
\begin{equation}
    8 R \sin \frac \alpha 3 \sin \frac \beta 3 \sin \frac \gamma 3,
\end{equation}
% Coxeter s. 23-25
% https://en.wikipedia.org/wiki/Hofstadter_points
a samo twierdzenie opiszą Coxeter \cite[s. 39-41]{coxeter_1967}, Zetel \cite[s. 177-179]{zetel_2020}, Neugebauer, Bogdańska \cite[s. 87, 105]{neugebauer_2018} (z dopiskiem, że nazwie się je kiedyś ,,ostatnim twierdzeniem geometrii elementarnej'' (?)). % lang-pl
e il teorema stesso è descritto da Coxeter \cite[p. 39-41]{coxeter_1967}, Zetel \cite[p. 177-179]{zetel_2020}, Neugebauer e Bogdańska \cite[p. 87, 105]{neugebauer_2018} (con l'osservazione che un giorno verrà chiamato ``l'ultimo teorema della geometria elementare'' (?)). % lang-it


Jak pamiętamy, pola powierzchnie wszystkich figur płaskich (trójkąta, równoległoboku, itd.) zostaną wyznaczone w starożytności przez podział figury na mniejsze części i ułożenie z nich prostokąta. % lang-pl
Come ricordiamo, le aree di tutte le figure piane (triangolo, parallelogramma ecc.) furono determinate nell'antichità suddividendole in parti più piccole e ricomponendole in un rettangolo. % lang-it

Próbowano tego samego wobec koła, bez powodzenia. % lang-pl
Si tentò di fare lo stesso con il cerchio, senza successo. % lang-it
\begin{problem}[quadratura del cerchio] % lang-it
\index{kwadratura koła}% % lang-it
    Costruire un quadrato avente area uguale a quella di un cerchio assegnato. % lang-it
\end{problem} % lang-it


\begin{problem}[kwadratura koła] % lang-pl
\index{kwadratura koła}% % lang-pl
    Wykreślić kwadrat o polu równym polu danego koła. % lang-pl
\end{problem} % lang-pl
Āpastamba, vissuto cinque secoli prima di Cristo e autore di uno dei quattro Śrautasūtra (testi vedici dedicati alla geometria), descrive come costruire un altare sacrificale. % lang-it

Apastamba, żyjący pięć wieków przed Chrystusem autor jednej z czterech śrautasutr (wedyjskiego tekstu poświęconemu geometrii), opisuje jak zbudować ołtarz ofiarny. % lang-pl
Oprócz próby kwadratury koła poda też przybliżenie $\sqrt 2$ z dokładnością do pięciu miejsc po przecinku. % lang-pl
Oltre a tentare la quadratura del cerchio, fornisce anche un'approssimazione di $\sqrt 2$ corretta a cinque cifre decimali. % lang-it


Opowiadaliśmy już o nie do końca trafionych pomysłach Mikołaja z Kuzy na stronie \pageref{nierownosc_mikolaja_z_kuzy}. % lang-pl
Abbiamo già parlato delle idee non del tutto corrette di Niccolò Cusano a pagina \pageref{nierownosc_mikolaja_z_kuzy}. % lang-it
\index[persons]{Mikołaj z Kuzy}% % lang-pl
\index[persons]{Niccolò Cusano}% % lang-it
% TODO: Nicolaus Cusanus, Nicolaus Krebs, Mikołaj Kuzańczyk
Adam Adamandy Kochański herbu Lubicz, dworzanin króla Jana Sobieskiego i członek zakonu jezuitów, znajdzie przybliżenie % lang-pl

Adam Adamandy Kochański dello stemma Lubicz, cortigiano del re Giovanni Sobieski e membro dell'ordine dei Gesuiti, troverà l'approssimazione % lang-it
\index[persons]{Kochański, Adam Adamandy}%
\begin{equation}
    \pi \approx \sqrt{\frac{40}{3} - 2 \sqrt{3}} = 3.14153\,33387...
\end{equation}
w pracy \emph{Observationes Cyclometricae ad facilitandam Praxin accommodatae} (Obserwacje cyklometryczne przystosowane dla ułatwienia praktycznego użycia) z 1685 roku. % lang-pl

Ludzie nie przestaną szukać innych niezbyt czasochłonnych przepisów. % lang-pl
nell'opera \emph{Observationes Cyclometricae ad facilitandam Praxin accommodatae} del 1685. % lang-it

Jacob de Gelder w 1849 znajdzie dokładniejsze przybliżenie % lang-pl
I matematici continueranno a cercare altre costruzioni semplici e pratiche. % lang-it

Nel 1849 Jacob de Gelder troverà un'approssimazione più accurata % lang-it
\index[persons]{Gelder@de Gelder, Jacob}%
\begin{equation}
    \pi \approx 3 + \frac{4^2}{7^2 + 8^2} = \frac{355}{113} = 3.14159\,29203...
\end{equation}

Ramanujan w 1914 roku poda inny przepis oparty o przybliżenie zgadzające się na ośmiu miejscach po przecinku: % lang-pl

Nel 1914 Ramanujan proporrà un'altra costruzione basata su un'approssimazione corretta a otto cifre decimali: % lang-it
\index[persons]{Ramanujan, Srinivasa}%
\begin{equation}
    \pi \approx \sqrt[4]{9^2 + \frac{19^2}{22}} = 3.14159\,26525...
\end{equation}

Konstrukcję Ramanujana przytoczy Jarosław Górnicki w $\Delta_{18}^3$. % lang-pl
La costruzione di Ramanujan è presentata da Jarosław Górnicki in $\Delta_{18}^3$. % lang-it


Ferdinand Lindemann odkryje w 1882 roku, że liczba $\pi$ pojawiająca się we wzorze na pole koła jest przestępna, zatem rozszerzenie $\mathbb Q(\sqrt{\pi}) / \mathbb Q$ jest za dużego stopnia i kwadratura koła, podobnie jak poprzednie dwa problemy, również nie jest wykonalna standardowymi narzędziami. % lang-pl
Nel 1882 Ferdinand Lindemann scoprirà che il numero $\pi$, che compare nella formula dell'area del cerchio, è trascendente; di conseguenza l'estensione $\mathbb Q(\sqrt{\pi})/\mathbb Q$ ha grado troppo elevato e la quadratura del cerchio, come i due problemi precedenti, non è realizzabile con gli strumenti classici. % lang-it

(Jeśli liczba $x$ jest konstruowalna, to stopień rozszerzenia $\mathbb Q(x) / \mathbb Q$ jest potęgą $2$, patrz \cite[s. 242]{hartshorne2000}; ale to nie wystarcza: wielomian $x^4 - x^3 - 5x^2 + 1$ ma cztery pierwiastki rzeczywiste $\alpha$ i rozszerzenia $\mathbb Q(\alpha)/\mathbb Q$ są stopnia $4$, ale grupa Galois ma rząd $12$ lub $24$, co nie jest potęgą dwójki). % lang-pl
(Se un numero $x$ è costruibile, allora il grado dell'estensione $\mathbb Q(x)/\mathbb Q$ è una potenza di $2$, vedi \cite[p. 242]{hartshorne2000}; ma questa condizione non è sufficiente: il polinomio $x^4-x^3-5x^2+1$ ha quattro radici reali $\alpha$ e le estensioni $\mathbb Q(\alpha)/\mathbb Q$ hanno grado $4$, mentre il gruppo di Galois ha ordine $12$ oppure $24$, che non sono potenze di due). % lang-it
\index[persons]{Lindemann, Ferdinand}

Kwadratura koła jest równoważna rektyfikacji okręgu, czyli zadaniu, by wykreślić odcinek tej samej długości, co obwód danego okręgu. % lang-pl

To również nie jest wykonalne. % lang-pl
La quadratura del cerchio è equivalente alla rettificazione del cerchio, cioè al problema di costruire un segmento avente la stessa lunghezza della circonferenza assegnata. % lang-it

Anche questo problema è impossibile. % lang-it

%
% eksperymentalne tłumaczenie na włoski

%

\section{Stożkowe} % lang-pl

\section{Coniche} % lang-it

% TODO: Alhazen => https://en.wikipedia.org/wiki/Straightedge_and_compass_construction => niemożliwy
% https://en.wikipedia.org/wiki/Alhazen%27s_problem

\begin{problem}[bilardowy Alhazena] % lang-pl
    Dany jest okrąg $\Gamma$ oraz dwa punkty $A$, $B$, obydwa w jego wnętrzu lub obydwa na zewnątrz. % lang-pl
\begin{problem}[problema biliardistico di Alhazen] % lang-it
    
    Skonstruować trójkąt równoramienny wpisany w okrąg $\Gamma$ tak, żeby punkt $A$ leżał na jednym ramieniu, zaś $B$ na drugim. % lang-pl
    Sono dati un cerchio $\Gamma$ e due punti $A$, $B$, entrambi al suo interno oppure entrambi al suo esterno. % lang-it
    
    Costruire un triangolo isoscele inscritto nel cerchio $\Gamma$ in modo che il punto $A$ appartenga a un lato obliquo e il punto $B$ all'altro. % lang-it
\end{problem}

Tematyką zajmie się już Ptolemeusz około II wieku naszej ery, ale dopiero XI-wieczny matematyk arabski, Alhazen (Hasan Ibn al-Haytham), sformułuje problem ogólniej, a oprócz tego przedstawi też jego rozwiązanie w książce \emph{Kitab al-Manazur} (łacińskie \emph{De aspectibus}, o optyce). % lang-pl

Około 1965 roku, amerykański matematyk Jack Elkin rozwiąże problem algebraicznie; pojawi się tam pierwiastek równania czwartego stopnia, więc trójkąta nie da się wykreślić samym cyrklem z linijką. % lang-pl
L'argomento sarà studiato già da Tolomeo intorno al II secolo d.C., ma soltanto il matematico arabo dell'XI secolo Alhazen (Hasan Ibn al-Haytham) ne formulerà una versione più generale e ne presenterà una soluzione nel \emph{Kitab al-Manazur} (in latino \emph{De aspectibus}), dedicato all'ottica. % lang-it

Al-Haytham i inni (na przykład Christiaan Huygens) wykorzystają wcześniej przecięcie stożkowych. % lang-pl
Intorno al 1965 il matematico statunitense Jack Elkin risolverà il problema algebricamente; comparirà una radice di un'equazione di quarto grado, per cui il triangolo non può essere costruito con la sola riga e il compasso. % lang-it
Al-Haytham e altri (per esempio Christiaan Huygens) utilizzeranno invece l'intersezione di coniche. % lang-it


Problem ma prostą interpretację fizyczną, dlatego też znalazł się w dziele o optyce. % lang-pl
Il problema ha una semplice interpretazione fisica, motivo per cui compare in un'opera di ottica. % lang-it

Promień światłą biegnie z punktu $A$ aż uderzy w powierzchnię kuli, po czym zmienia kierunek tak, że kąty między promieniem przed i po uderzeniu oraz normalną do powierzchni sfery są równe. % lang-pl
Un raggio luminoso parte dal punto $A$, colpisce la superficie di una sfera e cambia direzione in modo che gli angoli formati dal raggio incidente e riflesso con la normale alla superficie siano uguali. % lang-it

Szukamy takiego punktu, by po odbiciu trafić w punkt $B$. % lang-pl
Cerchiamo il punto di riflessione tale che il raggio raggiunga il punto $B$ dopo il rimbalzo. % lang-it


Więcej przeczytać można u Hartshorne'a \cite[s. 278]{hartshorne2000}. % lang-pl
Per ulteriori dettagli si veda Hartshorne \cite[p. 278]{hartshorne2000}. % lang-it

https://www.deltami.edu.pl/2017/06/zadanie-alhazena/

%
% eksperymentalne tłumaczenie na włoski

%

\section{Uszkodzone przyrządy} % lang-pl
\section{Strumenti danneggiati} % lang-it

Problemy od \ref{broken_ruler_compass_hartshorne_start} do \ref{broken_ruler_compass_hartshorne_end} pochodzą z książki Hartshorne'a \cite[s. 25, 26]{hartshorne2000}. % lang-pl
I problemi da \ref{broken_ruler_compass_hartshorne_start} a \ref{broken_ruler_compass_hartshorne_end} provengono dal libro di Hartshorne \cite[p. 25, 26]{hartshorne2000}. % lang-it

\begin{geoconstruction}
[połamana linijka] % lang-pl
[riga spezzata] % lang-it
\label{broken_ruler_compass_hartshorne_start}%
\index{linijka!połamana}%
    Dane są dwa punkty $A$ i $B$ na płaszczyźnie, odległe od siebie o około trzy nible. % lang-pl
    Mając do dyspozycji fragment linijki o długości jednej nibli oraz sprawny cyrkiel, narysować odcinek $AB$. % lang-pl
    Sono dati due punti $A$ e $B$ nel piano, distanti circa tre nibli. % lang-it
    Disponendo di un frammento di riga lungo una nible e di un compasso funzionante, tracciare il segmento $AB$. % lang-it
\end{geoconstruction}
% Hartshorne s. 25
O problemie pisze Eves \cite[s. 181]{eves1_1972}. % lang-pl
Del problema scrive Eves \cite[p. 181]{eves1_1972}. % lang-it

\begin{geoconstruction}
[zardzewiały cyrkiel] % lang-pl
[compasso arrugginito] % lang-it
    \index{cyrkiel!zardzewiały}
    Dane są dwa punkty $A$ i $B$ na płaszczyźnie, odległe od siebie o około pięć nibli. % lang-pl
    Mając do dyspozycji zardzewiały cyrkiel, którym można kreślić jedynie okręgi o promieniu dwóch nibli, skonstruować trójkąt równoboczny oparty o bok $AB$. % lang-pl
    Sono dati due punti $A$ e $B$ nel piano, distanti circa cinque nibli. % lang-it
    Disponendo di un compasso arrugginito con cui è possibile tracciare soltanto cerchi di raggio pari a due nibli, costruire un triangolo equilatero avente $AB$ come lato. % lang-it
\end{geoconstruction}

O zardzewiałym cyrklu będzie rozmyślać Abu al-Wafa Buzjani, matematyk perski, % lang-pl
Sul compasso arrugginito rifletterà Abu al-Wafa Buzjani, matematico persiano, % lang-it
\index[persons]{Buzjani, Abu al-Wafa}%
a pod koniec piętnastego wieku zastosują go w konstrukcjach Leonardo da Vinci oraz Albrecht Dürer. % lang-pl
e alla fine del XV secolo esso verrà impiegato nelle costruzioni di Leonardo da Vinci e Albrecht Dürer. % lang-it
\index[persons]{Dürer, Albrech}%
\index[persons]{da Vinci, Leonardo}%
% ŹRÓDŁO: https://en.wikipedia.org/wiki/Poncelet-Steiner_theorem#History

\begin{geoconstruction}
    Dany jest punkt $A$ leżący na prostej $l$. % lang-pl
    Skonstruować prostą prostopadłą do $l$ przechodzącą przez $A$ przy użyciu linijki i zardzewiałego cyrkla. % lang-pl
    È dato un punto $A$ appartenente alla retta $l$. % lang-it
    Costruire la retta perpendicolare a $l$ passante per $A$ usando una riga e un compasso arrugginito. % lang-it
\end{geoconstruction}
% Hartshorne s. 25

\begin{geoconstruction}
    Dany jest punkt $A$ leżący ponad cztery nible od prostej $l$. % lang-pl
    Skonstruować prostą prostopadłą do $l$ przechodzącą przez $A$ przy użyciu linijki i zardzewiałego cyrkla. % lang-pl
    È dato un punto $A$ a più di quattro nibli dalla retta $l$. % lang-it
    Costruire la retta perpendicolare a $l$ passante per $A$ usando una riga e un compasso arrugginito. % lang-it
\end{geoconstruction}
% Hartshorne s. 25

\begin{geoconstruction}
    Dane są trzy niewspółliniowe punkty $A$, $B$ oraz $C$. % lang-pl
    Skonstrować punkt $D$ na prostej $AC$ tak, żeby odcinki $AD$ oraz $AB$ były równej długości, przy użyciu linijki i zardzewiałego cyrkla. % lang-pl
    Sono dati tre punti non allineati $A$, $B$ e $C$. % lang-it
    Costruire un punto $D$ sulla retta $AC$ in modo che i segmenti $AD$ e $AB$ abbiano la stessa lunghezza, usando una riga e un compasso arrugginito. % lang-it
\end{geoconstruction}
% Hartshorne s. 26

\begin{geoconstruction}
    Dany jest odcinek $AB$ o długości ponad dwóch nibli oraz prosta $l$, która nie przechodzi przez końce odcinka. % lang-pl
    Skonstrować punkt $C$ na prostej $l$ tak, żeby odcinki $AB$ oraz $AC$ były równej długości, przy użyciu linijki i zardzewiałego cyrkla. % lang-pl
    Sono dati un segmento $AB$ di lunghezza superiore a due nibli e una retta $l$ che non passa per gli estremi del segmento. % lang-it
    Costruire un punto $C$ sulla retta $l$ in modo che i segmenti $AB$ e $AC$ abbiano la stessa lunghezza, usando una riga e un compasso arrugginito. % lang-it
\end{geoconstruction}
% Hartshorne s. 26

\begin{proposition}
\label{broken_ruler_compass_hartshorne_end}%
    Każdą konstrukcję geometryczną, którą można wykonać cyrklem i linijką, można powtórzyć przy użyciu zardzewiałego cyrkla i linijki. % lang-pl
    Qualsiasi costruzione geometrica realizzabile con riga e compasso può essere ripetuta usando una riga e un compasso arrugginito. % lang-it
\end{proposition}
% Hartshorne s. 26

Wynik ten zostanie odkryty relatywnie wcześnie. % lang-pl
Pierwszy pokaże go Lodovico Ferrari, uczeń Gerolamo Cardano i odkrywca metody rozwiązywania równań czwartego stopnia w pojedynku przeciwko Niccolò Fontanie Tartaglii. % lang-pl
Questo risultato sarà scoperto relativamente presto. % lang-it
Il primo a dimostrarlo sarà Lodovico Ferrari, allievo di Gerolamo Cardano e scopritore del metodo per risolvere le equazioni di quarto grado nella celebre disputa con Niccolò Fontana Tartaglia. % lang-it
\index[persons]{Ferrari, Lodovico}%
\index[persons]{Cardano, Gerolamo}%
W przeciągu dziesięciu lat wyczyn ten powtórzą Cardano, Tartaglia oraz uczeń Tartaglii, Benedetti. % lang-pl
Nel giro di dieci anni il risultato sarà ottenuto nuovamente da Cardano, Tartaglia e dall'allievo di Tartaglia, Benedetti. % lang-it
\index[persons]{Tartaglia, Nicolo}%
\index[persons]{Benedetti, Giambattista}%

Ciąg dalszy to oczywiście twierdzenie \ref{thm:poncelet_steiner}, że zardzewiałego cyrkla wystarczy użyć raz. % lang-pl
Il seguito naturale è naturalmente il teorema \ref{thm:poncelet_steiner}, secondo cui è sufficiente usare il compasso arrugginito una sola volta. % lang-it

%
% eksperymentalne tłumaczenie na włoski

%

\section{Konstrukcje bez linijki} % lang-pl

\section{Costruzioni senza riga} % lang-it

% https://www.deltami.edu.pl/2015/08/prosto-w-srodek/ % TODO

Coxeter \cite[s. 95]{coxeter_1967} sugeruje skonstruować (samym) cyrklem wierzchołki sześciokąta foremnego, odcinek dwa razy dłuższy (i krótszy) od danego, obraz inwersyjny danego punktu leżącego dowolnie blisko środka i wreszcie podzielić dany odcinek na $n$ równych części. % lang-pl
Coxeter \cite[p. 95]{coxeter_1967} suggerisce di costruire con il solo compasso i vertici di un esagono regolare, un segmento lungo il doppio (e la metà) di un segmento dato, l'immagine inversiva di un punto dato situato arbitrariamente vicino al centro e infine di dividere un segmento dato in $n$ parti uguali. % lang-it


Okazuje się, że do żadnej klasycznej konstrukcji nie jest potrzebna linijka. % lang-pl
Si scopre che nessuna costruzione classica richiede realmente l'uso della riga. % lang-it

\begin{theorem}[twierdzenie Mohra-Mascheroniego] % lang-pl
\begin{theorem}[teorema di Mohr-Mascheroni] % lang-it
\index{twierdzenie!Mohra-Mascheroniego}%
    Jeśli dana konstrukcja jest wykonalna za pomocą cyrkla i~linijki, to jest ona wykonalna za pomocą samego cyrkla, o ile pominiemy rysowanie linii i ograniczymy się do wyznaczania punktów konstrukcji. % lang-pl
    
    Se una costruzione è realizzabile con riga e compasso, allora è realizzabile con il solo compasso, purché si ometta il tracciamento delle linee e ci si limiti a determinare i punti della costruzione. % lang-it
\end{theorem}
% PRZECZYTANO: https://en.wikipedia.org/wiki/Mohr-Mascheroni_theorem

Wynik będzie mieć ciekawą historię. % lang-pl

Pierwszy znajdzie go Georg Mohr, włoski geometra i~poeta w \emph{Euclides Danicus} (1672), ale jego odkrycie popadnie w zapomnienie, po czym będzie tak marnieć aż do roku 1928! % lang-pl
Questo risultato avrà una storia interessante. % lang-it

A scoprirlo per primo sarà Georg Mohr, geometra e poeta danese, nell'opera \emph{Euclides Danicus} (1672), ma la sua scoperta cadrà nell'oblio e vi rimarrà fino al 1928! % lang-it
% established his results by using the idea of reflection in a line.
% Euclides Danicus = duński Euklides
\index[persons]{Mohr, Georg}%
Niezależnie twierdzenie odkryje Lorenzo Mascheroni; wykorzysta odbicia względem prostych i~opisze je w~\emph{La Geometria del Compasso} (1797). % lang-pl

Indipendentemente, il teorema sarà riscoperto da Lorenzo Mascheroni, che utilizzerà le riflessioni rispetto alle rette e le descriverà nella sua opera \emph{La Geometria del Compasso} (1797). % lang-it
\index[persons]{Mascheroni, Lorenzo}%
W 1890 jeszcze raz to samo, ale nie tak samo (inwersjami) zrobi śląski\footnote{Urodzony w Księstwie Górnego i Dolnego Śląska, znanym także jako Śląsk Austriacki} matematyk August Adler w pracy, której tytuł nadgryzie ząb czasu. % lang-pl

Nel 1890 il matematico slesiano\footnote{Nato nel Ducato dell'Alta e Bassa Slesia, noto anche come Slesia austriaca.} August Adler otterrà nuovamente lo stesso risultato, ma con un metodo diverso, basato sulle inversioni, in un lavoro il cui titolo è stato ormai consumato dal tempo. % lang-it
% August Adler (24 January 1863, Opava, Austrian Silesia – 17 October 1923, Vienna) was a Czech and Austrian mathematician noted for using the theory of inversion to provide an alternate proof of Mascheroni's compass and straightedge construction theorem. Czyli nie był z niczego innego znany?
Wreszcie w 1928 student Hjelmsleva przeglądając księgarnię w Kopenhadze i jej zawartość znajdzie książkę Mohra. % lang-pl

Infine, nel 1928, uno studente di Hjelmslev troverà il libro di Mohr sfogliando i volumi di una libreria di Copenaghen. % lang-it
% źródło tej rewelacji: George Martin - Geometric Constructions (1998), s. 54
Będzie to duża niespodzianka dla Hjelmsleva. % lang-pl
Sarà una grande sorpresa per Hjelmslev. % lang-it


Dla dowodu wystarczy pokazać, że cyrkel wystarcza do przeprowadzenia pięciu konstrukcji: % lang-pl
Per la dimostrazione basta mostrare che il compasso è sufficiente per eseguire le seguenti cinque costruzioni: % lang-it
\begin{enumerate}
    \item poprowadzenia prostej przez dwa punkty, % lang-pl
    \item wykreślenia okręgu o danym środku i promieniu, % lang-pl
    \item znalezienia punktu przecięcia dwóch nierównoległych prostych, % lang-pl
    \item znalezienia punktu (lub punktów) przecięcia prostej z okręgiem, % lang-pl
    \item znalezienia punktu (lub punktów) przecięcia dwóch okręgów. % lang-pl
    \item tracciare la retta passante per due punti, % lang-it
    \item tracciare un cerchio di centro e raggio assegnati, % lang-it
    \item trovare il punto di intersezione di due rette non parallele, % lang-it
    \item trovare il punto (o i punti) di intersezione tra una retta e un cerchio, % lang-it
    \item trovare il punto (o i punti) di intersezione di due cerchi. % lang-it
\end{enumerate}

O wyniku Mohra-Mascheroniego napisze Eves \cite[s. 169-172]{eves1_1972}. % lang-pl

Del risultato di Mohr-Mascheroni scrive Eves \cite[p. 169-172]{eves1_1972}. % lang-it

%

\begin{problem}[zadanie Napoleona] % lang-pl
\begin{problem}[problema di Napoleone] % lang-it
\index{zadanie!Napoleona}
\label{napoleon_problem}
	Dany jest okrąg oraz jego środek. % lang-pl
	
	Podzielić okrąg na cztery łuki równej miary korzystając z cyrkla, ale nie linijki. % lang-pl
	Sono dati un cerchio e il suo centro. % lang-it
	
	Dividere il cerchio in quattro archi di uguale ampiezza usando il compasso ma non la riga. % lang-it
\end{problem}
% PRZECZYTANO: https://en.wikipedia.org/wiki/Napoleon%27s_problem

Jak zawsze w przypadku Napoleona i geometrii, nie wiadomo, czy rozwiązał albo choć wymyślił przedstawione wyżej zadanie konstrukcyjne. % lang-pl

Come spesso accade quando si parla di Napoleone e geometria, non si sa se abbia risolto o persino formulato il problema di costruzione sopra descritto. % lang-it
\index[persons]{Bonaparte, Napoleon}%
W trudniejszej wersji okrąg nie ma środka i trzeba go najpierw znaleźć. % lang-pl

Rozwiązanie podają Bogdańska, Neugebauer \cite[s. 116]{neugebauer_2018} (z wykorzystaniem okręgów Torricelliego); % lang-pl
Nella versione più difficile il cerchio non ha il centro indicato e bisogna trovarlo per prima cosa. % lang-it

u Audina \cite[s. 105]{audin_2003} jest to ćwiczenie. % lang-pl
Una soluzione è fornita da Bogdańska e Neugebauer \cite[p. 116]{neugebauer_2018} (utilizzando i cerchi di Torricelli); % lang-it

in Audin \cite[p. 105]{audin_2003} compare invece come esercizio. % lang-it
\index{okrąg Torricelliego}%

%

% eksperymentalne tłumaczenie na włoski

```latex
%

\section{Konstrukcje bez cyrkla} % lang-pl

\section{Costruzioni senza compasso} % lang-it

\begin{theorem}[twierdzenie Ponceleta-Steinera] % lang-pl
\begin{theorem}[teorema di Poncelet-Steiner] % lang-it
\label{thm:poncelet_steiner}%
\index{twierdzenie!Ponceleta-Steinera}%
    Jeśli dana konstrukcja jest wykonalna za pomocą cyrkla i linijki, to jest ona wykonalna za pomocą samej linijki, o ile dany jest na płaszczyźnie pewien okrąg wraz ze środkiem. % lang-pl

    Se una costruzione è realizzabile con riga e compasso, allora è realizzabile con la sola riga, purché nel piano sia dato un cerchio insieme al suo centro. % lang-it
    \index{twierdzenie!Ponceleta-Steinera}%
\end{theorem}

Nazwa twierdzenia pochodzi od dwóch panów: Jeana Ponceleta, który postawi w 1822 roku hipotezę oraz Jakoba Steinera, który dostarczy brakujący dowód w 1833 roku. % lang-pl

Il nome del teorema deriva da due studiosi: Jean Poncelet, che formulò la congettura nel 1822, e Jakob Steiner, che fornì la dimostrazione mancante nel 1833. % lang-it
\index[persons]{Poncelet, Jean}%
\index[persons]{Steiner, Jakob}%
Zamiast całego okręgu, wystarczy, że mamy do dyspozycji mały jego łuk -- to wzmocnienie, o które pokusi się Francesco Severi w 1904 roku. % lang-pl

Al posto dell'intero cerchio, è sufficiente avere a disposizione un suo piccolo arco: questo rafforzamento fu ottenuto da Francesco Severi nel 1904. % lang-it
\index[persons]{Severi, Francesco}%

Sama linijka to za mało, bo nie pozwala na wyciąganie pierwiastków kwadratowych, tak jak opisze to Łukasz Rajkowski w $\Delta_{15}^8$. % lang-pl

La sola riga non basta, poiché non permette di estrarre radici quadrate, come descrive Łukasz Rajkowski in $\Delta_{15}^8$. % lang-it
\index[persons]{Rajkowski, Łukasz}%
Ale zamiast środka okręgu możemy wymagać: % lang-pl

Tuttavia, al posto del centro del cerchio possiamo richiedere: % lang-it
\begin{itemize}
    \item równoległoboku, % lang-pl
    \item drugiego koncentrycznego okręgu, % lang-pl
    \item drugiego okręgu przecinającego pierwszy, % lang-pl
    \item drugiego okręgu rozłącznego z pierwszym i punktu na prostej łączącej ich środki albo ich osi potęgowej... % lang-pl
    \item un parallelogramma, % lang-it
    \item un secondo cerchio concentrico, % lang-it
    \item un secondo cerchio che interseca il primo, % lang-it
    \item un secondo cerchio disgiunto dal primo e un punto sulla retta che congiunge i loro centri oppure sul loro asse radicale... % lang-it
\end{itemize}

Zapewne istnieje jeszcze więcej takich warunków; o części napisze Eves \cite[s. 174-182]{eves1_1972}. % lang-pl

August Adler wspomniany we wcześniejszej sekcji zauważy, że linijka z dwoma brzegami albo ekierka pozwalają wykonać wszystkie klasyczne konstrukcje. % lang-pl
Esistono probabilmente molte altre condizioni di questo tipo; alcune sono descritte da Eves \cite[p. 174-182]{eves1_1972}. % lang-it
August Adler, menzionato nella sezione precedente, osservò che una riga a doppio bordo oppure una squadra consentono di eseguire tutte le costruzioni classiche. % lang-it


Dwie trójkątne linijki (w kształcie litery L) pozwalają podwoić sześcian (problem \ref{problem_delijski}) i dzielić kąt na trzy równe (problem \ref{trysekcja_kata}). % lang-pl
Due righe triangolari (a forma di L) permettono di duplicare il cubo (problema \ref{problem_delijski}) e di trisecare un angolo (problema \ref{trysekcja_kata}). % lang-it

\begin{problem}
    Dany jest okrąg $\Gamma$ ze środkiem $O$ oraz odcinek. % lang-pl
    
    Skonstruować środek odcinka. % lang-pl
    Sono dati un cerchio $\Gamma$ con centro $O$ e un segmento. % lang-it
    
    \hfill \emph{(15 kroków)}. % Hartshorne s. 192 % lang-pl
    Costruire il punto medio del segmento. % lang-it
    \hfill \emph{(15 passi)}. % Hartshorne s. 192 % lang-it
\end{problem}

\begin{problem}
    Dany jest okrąg $\Gamma$ ze środkiem $O$, prosta $l$ oraz punkt $P$. % lang-pl
    
    Skonstruować prostą równoległą do $l$, która przechodzi przez $P$. % lang-pl
    Sono dati un cerchio $\Gamma$ con centro $O$, una retta $l$ e un punto $P$. % lang-it
    
    \hfill \emph{(16 kroków)}. % Hartshorne s. 192 % lang-pl
    Costruire la retta parallela a $l$ passante per $P$. % lang-it
    \hfill \emph{(16 passi)}. % Hartshorne s. 192 % lang-it
\end{problem}

\begin{problem}
    Dany jest okrąg $\Gamma$ ze środkiem $O$, odcinek $OA$ oraz półprosta $OB$. % lang-pl
    
    Skonstruować punkt $C$ na półprostej taki, że $OA$ i $OC$ są przystające. % lang-pl
    Sono dati un cerchio $\Gamma$ con centro $O$, il segmento $OA$ e la semiretta $OB$. % lang-it
    
    \hfill \emph{(17 kroków)}. % Hartshorne s. 193 % lang-pl
    Costruire un punto $C$ sulla semiretta tale che $OA$ e $OC$ siano congruenti. % lang-it
    \hfill \emph{(17 passi)}. % Hartshorne s. 193 % lang-it
\end{problem}

\begin{problem}
    Dany jest okrąg $\Gamma$ ze środkiem $O$, prosta $l$ oraz punkt $P$. % lang-pl
    
    Skonstruować prostą prostopadłą do $l$, która przechodzi przez $P$. % lang-pl
    Sono dati un cerchio $\Gamma$ con centro $O$, una retta $l$ e un punto $P$. % lang-it
    
    \hfill \emph{(33 kroki)}. % Hartshorne s. 193 % lang-pl
    Costruire la retta perpendicolare a $l$ passante per $P$. % lang-it
    \hfill \emph{(33 passi)}. % Hartshorne s. 193 % lang-it
\end{problem}

Zetel \cite[s. 17]{zetel_2020} poda wariację tego problemu: jego prosta $l$ jest średnicą okręgu. % lang-pl

Premier Rosji Michaił Miszustin odwiedzi we wrześniu 2021 roku liceum imienia P. L. Kapicy w Dołgoprudnym i~poprosi uczniów o~rozwiązanie jeszcze prostszego zadania, gdzie punkt $P$ leży na okręgu $\Gamma$. % lang-pl
 % https://www.theguardian.com/science/2021/sep/20/did-you-solve-it-russias-prime-minister-sets-a-geometry-puzzle % lang-pl
Zetel \cite[p. 17]{zetel_2020} propone una variante di questo problema: la sua retta $l$ è un diametro del cerchio. % lang-it
Nel settembre 2021 il primo ministro russo Michail Mišustin visiterà il liceo P. L. Kapica di Dolgoprudnyj e chiederà agli studenti di risolvere un problema ancora più semplice, nel quale il punto $P$ giace sul cerchio $\Gamma$. % https://www.theguardian.com/science/2021/sep/20/did-you-solve-it-russias-prime-minister-sets-a-geometry-puzzle % lang-it

\begin{problem}
    Dany jest okrąg $\Gamma$ ze środkiem $O$, prosta $l$ oraz dwa punkty $A$, $B$. % lang-pl
    
    Skonstruować punkt, gdzie prosta $l$ przecina okrąg o środku $A$ i promieniu $AB$. % lang-pl
    Sono dati un cerchio $\Gamma$ con centro $O$, una retta $l$ e due punti $A$, $B$. % lang-it
    
    \hfill \emph{(54 kroki)}. % Hartshorne s. 193 % lang-pl
    Costruire il punto in cui la retta $l$ interseca il cerchio di centro $A$ e raggio $AB$. % lang-it
    \hfill \emph{(54 passi)}. % Hartshorne s. 193 % lang-it
\end{problem}

\begin{problem}
    Dany jest okrąg $\Gamma$ ze środkiem $O$ oraz punkt $A$ poza okręgiem. % lang-pl
    
    Skonstruować styczne do $\Gamma$ przechodzące przez $A$. % lang-pl
    Sono dati un cerchio $\Gamma$ con centro $O$ e un punto $A$ esterno al cerchio. % lang-it
    
    Costruire le tangenti a $\Gamma$ passanti per $A$. % lang-it
\end{problem}

\begin{problem}
    Dany jest okrąg $\Gamma$ ze środkiem $O$. % lang-pl
    
    Skonstruować pięciokąt foremny wpisany w $\Gamma$. % lang-pl
    È dato un cerchio $\Gamma$ con centro $O$. % lang-it
    
    \hfill \emph{(około 50 kroków)}. % Hartshorne s. 193 % lang-pl
    Costruire un pentagono regolare inscritto in $\Gamma$. % lang-it
    \hfill \emph{(circa 50 passi)}. % Hartshorne s. 193 % lang-it
\end{problem}

https://www.deltami.edu.pl/2017/06/z-sama-linijka-na-okrag-styczna-do-okregu-brakujacy-punkt/

https://www.deltami.edu.pl/2014/02/sama-linijka-mozna-nakreslic-okrag/

%
```

% eksperymentalne tłumaczenie na włoski

%

Konstrukcje od \ref{delta_2024_12_start} do \ref{delta_2024_12_end} opisze Stanisław Majchrzak w $\Delta_{24}^{12}$. % lang-pl
Le costruzioni da \ref{delta_2024_12_start} a \ref{delta_2024_12_end} sono descritte da Stanisław Majchrzak in $\Delta_{24}^{12}$. % lang-it

\begin{geoconstruction}
\label{delta_2024_12_start}%
    Dane jest pięć punktów leżących na okręgu $\Gamma$. % lang-pl
    Wykreślić styczną do $\Gamma$ w jednym z tych punktów. % lang-pl
    %
    Sono dati cinque punti appartenenti a un cerchio $\Gamma$. % lang-it
    Costruire la tangente a $\Gamma$ in uno di questi punti. % lang-it
\end{geoconstruction}

Uzasadnienie poprawności tej konstrukcji wymaga znajomości twierdzenia Pascala, patrz $\Delta_{14}^9$. % lang-pl
\index{twierdzenie!Pascala}% % lang-pl
La dimostrazione della correttezza di questa costruzione richiede la conoscenza del teorema di Pascal; si veda $\Delta_{14}^9$. % lang-it
\index{teorema!di Pascal} % lang-it

\begin{geoconstruction}
    Dane jest pięć punktów leżących na okręgu $\Gamma$ oraz prosta $l$ przechodząca przez jeden z nich. % lang-pl
    Wyznaczyć drugi punkt przecięcia okręgu $\Gamma$ oraz prostej $l$. % lang-pl
    %
    Sono dati cinque punti appartenenti a un cerchio $\Gamma$ e una retta $l$ passante per uno di essi. % lang-it
    Determinare il secondo punto d'intersezione tra il cerchio $\Gamma$ e la retta $l$. % lang-it
\end{geoconstruction}

Rozwiązanie zostanie wydrukowane w $\Delta_{17}^6$. % lang-pl
La soluzione sarà pubblicata in $\Delta_{17}^6$. % lang-it

\begin{geoconstruction}
    Dane są dwa okręgi przecinające się w dwóch punktach. % lang-pl
    Wyznaczyć środek jednego z nich. % lang-pl
    %
    Sono dati due cerchi che si intersecano in due punti. % lang-it
    Determinare il centro di uno di essi. % lang-it
\end{geoconstruction}

W kolejnych konstrukcjach przydatna jest znajomość twierdzenia: biegunowe dowolnego punktu $P$ względem okręgów należących do jednego pęku są współpękowe. % lang-pl
Nelle costruzioni successive è utile conoscere il seguente teorema: le polari di un punto arbitrario $P$ rispetto ai cerchi appartenenti a uno stesso fascio sono concorrenti. % lang-it

\begin{geoconstruction}
    Dane są dwa okręgi oraz punkt $A$ na zewnątrz jednego z nich (wewnątrz obydwu). % lang-pl
    Wykreślić okrąg przechodzący przez $A$, który jest współpękowy z dwoma danymi. % lang-pl
    %
    Sono dati due cerchi e un punto $A$ esterno a uno di essi (e interno a entrambi). % lang-it
    Costruire un cerchio passante per $A$ che appartenga allo stesso fascio dei due cerchi assegnati. % lang-it
\end{geoconstruction}

\begin{geoconstruction}
\label{delta_2024_12_end}%
    Dane są cztery (w trudniejszej wersji: trzy) okręgi takie, że żadne trzy nie są współpękowe. % lang-pl
    Wyznaczyć środek przynajmniej jednego z nich. % lang-pl
    %
    Sono dati quattro cerchi (nella versione più difficile: tre) tali che nessuna terna appartenga allo stesso fascio. % lang-it
    Determinare il centro di almeno uno di essi. % lang-it
\end{geoconstruction}

Dla rozłącznych okręgów należących do jednego pęku, środka nie da się (!) skonstruować. % lang-pl
Per cerchi disgiunti appartenenti a uno stesso fascio, il centro non può (!) essere costruito. % lang-it

%
% eksperymentalne tłumaczenie na włoski

%

\section{Przyrządy potężniejsze niż cyrkiel i linjka bez podziałki} % lang-pl
\section{Strumenti più potenti del compasso e della riga non graduata} % lang-it

\subsection{Konstrukcja neusis} % lang-pl
\subsection{Costruzione per neusis} % lang-it

Konstrukcja neusis z linijką z podziałką. % lang-pl
Costruzione per neusis mediante una riga graduata. % lang-it
%  νεύειν (neuein) 'incline towards'; plural: νεύσεις, neuseis => https://en.wikipedia.org/wiki/Neusis_construction
% https://en.wikipedia.org/wiki/Straightedge_and_compass_construction#Markable_rulers
Archimedes poda konstrukcję neusis dla siedmiokąta foremnego. % lang-pl
Archimede descriverà una costruzione per neusis del settagono regolare. % lang-it

\subsection{Konstrukcja stożkowa} % lang-pl
\subsection{Costruzioni coniche} % lang-it

Oprócz cyrkla i linijki mamy do dyspozycji narzędzie, które kreśli stożkowe o zadanych ognisku, kierownicy i mimośrodzie. % lang-pl
Taką samą moc ma zestaw złożony z cyrkla, linijki i wykresu paraboli $y = x^2$ dla $x \in [0, 1]$ albo sznurka, który pozwala kreślić elipsy o zadanej półosi i ogniskach. % lang-pl
Starożytni Grecy będą wiedzieć, że podwojenie sześcianu oraz trysekcja dowolnego kąta mają konstrukcje stożkowe. % lang-pl
Kwadratura koła nie ma. % lang-pl
Oltre al compasso e alla riga disponiamo di uno strumento capace di tracciare coniche assegnandone fuoco, direttrice ed eccentricità. % lang-it
La stessa potenza costruttiva si ottiene con un compasso, una riga e il grafico della parabola $y=x^2$ per $x \in [0,1]$, oppure con una cordicella che permetta di tracciare ellissi con semiassi e fuochi assegnati. % lang-it
Gli antichi Greci sapevano che la duplicazione del cubo e la trisezione di un angolo arbitrario ammettono costruzioni coniche. % lang-it
La quadratura del cerchio no. % lang-it
% \todofoot{Solid constructions} % https://en.wikipedia.org/wiki/Straightedge_and_compass_construction#Solid_constructions % https://en.wikipedia.org/wiki/Straightedge_and_compass_construction#Angle_trisection_2

\subsection{Origami} % lang-pl
\subsection{Origami} % lang-it

Patrz sekcja \ref{section_origami}. % lang-pl
Si veda la sezione \ref{section_origami}. % lang-it

%

%